% ============================================================
%  Harald Høffding, OM GRUNDLAGET FOR DEN HUMANE ETHIK.
%  "On the Foundation of Humane Ethics"
%  Kjøbenhavn: Andr. Fred. Høst & Søns Forlag, 1876 (Sally B. Salomons
%  Tryk).  [IV], 146 pp.
%
%  TRANSLATION (English), from transcription.tex (Danish) in this directory,
%  the text of record (transcribed and collated 2026-10-05; see its header).
%  Translated from the transcription as committed at 5af2fd8, 2026-10-05, then
%  brought into line with the repairs below (same day).
%
%  DRAFT COMPLETE (2026-10-05): the whole book, pp. 1--146, translated in reading
%  order; \opage 1--146 (p. 6 blank in the print) in order, once each; 44 notes,
%  as in the Danish, page by page.  Pp. 1--5 translated in the main session,
%  pp. 7--146 in eight parts by subagents working to this header's conventions and
%  TERMS, then spliced and harmonized.  To be read by an independent reviewer.
%  TRANSCRIPTION REPAIRED WHILE TRANSLATING (2026-10-05, TRANSLATION-PLAYBOOK.md
%  §6c): the 59 odd spots the translation found were read at the image and 27
%  corrections made through patches.POST (see the transcription's header); the
%  English follows the repaired text (false paragraph breaks pp. 44, 66, 103
%  joined; Bonald's quotation closed p. 64; letterspaced names; couplets pp. 82,
%  122 as one block; misprint comments pp. 24, 28, 49, 61, 143).
%  Richard III (p. 140) is given in Shakespeare's English (I, 1), with
%  Høffding's Danish in a comment.
%
%  BOOK-SPECIFIC CONVENTIONS (beyond TRANSLATION-PLAYBOOK.md §2):
%   * Title: "humane Ethik" = humane ethics, as in the repo's other Høffding
%     translations (Etik, Psykologi), which cite this book as "On the
%     Foundation of Humane Ethics".  Høffding opposes it to Christian ethics.
%   * Six chapters, unnumbered in the print, each under a centred title
%     (\chap, as in the transcription).  The fifth chapter's head is printed
%     with a leading dash, kept.
%   * Letterspacing in the Danish (\emph: names, stressed words) -> \emph;
%     italics (\textit) -> \textit.
%   * Footnotes per page as printed, \footnote[k]{...} at the same anchor,
%     marked *), **), ***) as in the print.  "Smlgn." / "Kfr." / "Jfr." ->
%     "Cf."; "Se" -> "See"; "Overs." -> "trans."; page and volume numbers,
%     "S." and "p.", as printed.
%   * Quotations Høffding gives in Danish translation are translated from
%     his Danish; what he quotes in another language stays as printed.
%     Danish verse is translated in the text, the Danish in a comment below.
%     Titles cited stay as cited (Høffding's own works in Danish).
%   * »...«, «...» and „...“ alike -> ``...''.
%   * "d. v. s." and "ɔ:" -> "i.e."; "f. Ex." -> "e.g."
%   * Misprints (PRINTED AS IS in the Danish): translate the sense, comment.
%   * Page markers \opage{N} inline at the English word where printed page N
%     begins (p. 6 is blank in the print).
%
%  TERMS (aligned with the repo's Etik and Psykologi translations):
%    Ethik / Ethiker / det Ethiske = ethics / ethicist / the ethical;
%    sædelig / Sædelighed = moral / morality; Pligt = duty;
%    Handlen = action (acting); Handling = action, act;
%    Formaal = end; Grundlag = foundation; Begrundelse = grounding;
%    Maalestok = standard; Sindelag = disposition; Drift = impulse;
%    Lyksalighed(sdriften) = (the impulse toward) happiness [felicity];
%    Lykke = happiness; Velfærd = welfare; Nytte = utility;
%    Lyst / Ulyst / Smerte = pleasure / displeasure / pain;
%    Følelse = feeling; Fornemmelse = sensation; Forestilling = idea;
%    Erkjendelse = knowledge (cognition where the act is meant);
%    Anerkjendelse = recognition; Bevidsthed = consciousness;
%    Villie = will; Villiens Frihed = freedom of the will;
%    Tilregnelighed = accountability; Avtoritet = authority;
%    Myndighed = authority, power (by context; "Straffemyndighed" = penal
%    authority); Ærefrygt = reverence; Samvittighed = conscience;
%    Opdragelse = education; Slægten = the race; Menneskeslægten = the human
%    race; Samfund = society; Fremskridt = progress; Udvikling = development
%    (evolution where the Darwinian theory is meant); Trin = stage;
%    Livsanskuelse = view of life; Uegennyttighed = disinterestedness;
%    Selvopholdelse = self-preservation; Kampen for Tilværelsen = the
%    struggle for existence; Fornuft = reason; Aand / aandelig = spirit /
%    spiritual (mental where plainly psychological); sjælelig = mental.
%  Settled with pp. 7--18: Erkjendelse = cognition where set against
%    feeling or action, knowledge elsewhere; Anlæg = predisposition;
%    Evne = capacity; Stemning = mood; Sansning = sensation; Øjemed =
%    purpose; Villen = willing; Tilvanthed = habituation; Sædvane = custom;
%    vilkaarlig = voluntary (uvilkaarlig = involuntary; arbitrary where
%    that is the sense).
%  Settled with pp. 19--107: Forpligtelse = obligation; Nyttemoral /
%    Sympathimoral / Pligtmoral = morality of utility / of sympathy / of
%    duty; Udviklingshypothesen = the hypothesis of evolution;
%    Kvalitetsvalg = natural selection; Kampen for Livet, Livskampen = the
%    struggle for life; Forbillede = model; Menneskekjærlighed = love of
%    humanity; Humanitetens Rige = the kingdom of humanity; Slægt =
%    generation where generations are contrasted, the race elsewhere;
%    Underkastelse = submission; Lydighed = obedience; Hengivenhed =
%    devotion; Hengivelse = self-surrender; Pietet = piety; Grunddyd =
%    fundamental virtue; Tro = faith (Christian), belief (in spirits,
%    gods); Skik = custom; Vedtægt = convention; Klogskab = prudence;
%    Afhængighedsfølelse = feeling of dependence; Aandsliv = spiritual
%    life; Magt / Ret (as stages) = might / right; Fremadskriden = advance.
%    Names as printed (Leibnitz), classical names in English (Plato).
% ============================================================
\documentclass[12pt,a4paper,openany]{book}
\usepackage[utf8]{inputenc}
\usepackage[T1]{fontenc}
\usepackage{libertinus}
\usepackage{libertinust1math}
\usepackage{textalpha}
\usepackage[english]{babel}
\usepackage[protrusion=true,expansion=true]{microtype}
\usepackage{enumitem}
\usepackage{setspace}
\usepackage[margin=1.2in]{geometry}
\usepackage{fancyhdr}
\setlength{\headheight}{28pt}
\usepackage{hyperref}
\hypersetup{pdftitle={On the Foundation of Humane Ethics (1876)}, pdfauthor={Harald Høffding}, hidelinks}
\pagestyle{fancy}
\fancyhf{}
\fancyhead[LE,RO]{\thepage}
\fancyhead[RE]{\textit{Harald Høffding: On the Foundation of Humane Ethics}}
\fancyhead[LO]{\textit{1876}}
\fancypagestyle{plain}{\fancyhf{}\fancyhead[LE,RO]{\thepage}}
\setstretch{1.3}
\usepackage{marginnote}
\renewcommand{\thefootnote}{%
  \ifcase\value{footnote}\or *)\or **)\or ***)\else \arabic{footnote})\fi}
\newcommand{\tocline}[2]{\noindent #1\dotfill\ #2\par}
\newcommand{\opage}[1]{\marginnote{\footnotesize\textit{[#1]}}}
% a chapter of the print: new page, centred title, a bookmark and a contents line
\newcommand{\chap}[1]{\clearpage\addcontentsline{toc}{chapter}{#1}\vspace*{2em}\begin{center}{\large #1}\end{center}\medskip}
\pdfstringdefDisableCommands{\renewcommand{\footnote}[2][]{}\renewcommand{\opage}[1]{}}

\begin{document}

% --- title page ---
\begin{titlepage}
\begin{center}
\vspace*{3em}
{\large ON THE FOUNDATION}\\[0.8em]
OF\\[0.8em]
{\LARGE HUMANE ETHICS}\\[2em]
BY\\[1em]
DR. HARALD HØFFDING\\[2em]
\vfill
{\small Translated from the Danish}\\[0.6em]
{\small \textit{Om Grundlaget for den humane Ethik.}\\
Kjøbenhavn: Andr. Fred. Høst \& Søns Forlag, 1876}
\end{center}
\end{titlepage}
% --- verso: the motto (German, as printed) ---
\thispagestyle{empty}
\vspace*{2em}
\begin{flushright}
Was ist das Heiligste? --- Das, was heut und ewig die Geister,\\
Tiefer und tiefer gefühlt, immer nur einiger macht.\\
\hfill Goethe.
\end{flushright}
\clearpage
% --- Contents ---
\thispagestyle{empty}
\begin{center}{\large \emph{Contents.}}\end{center}
\hfill Page\par
\tocline{Introduction}{1.}
\tocline{Psychological Remarks}{7.}
\tocline{Individual Starting Points}{19.}
\tocline{Authority}{55.}
\tocline{The Ethical Law and Progress}{87.}
\tocline{Freedom of the Will}{108.}
\medskip
\noindent{\small [Page numbers are those of the Danish edition, given in the margin of this translation.]}
\begin{center}\rule{4em}{0.4pt}\end{center}
\clearpage

% ===== INTRODUCTION (pp. 1--5) =====
\chap{Introduction.}
\opage{1}Although the great struggles of recent times between the various
religious, social and political tendencies also concern ethical questions,
one may nevertheless claim that, as regards the content of ethics, the law
of duty itself, there prevails far greater uniformity and agreement than
the opponents in the heat of the dispute are aware of. At the same stage of
culture the general principles and rules of human action will, by the
course of things, be the same for all. A common ground of humanity makes up
the arena in which the struggle is fought. From this standpoint even the
opposition between Christian and humane ethics disappears. ``The
Christian,'' says \emph{Martensen}, ``is in our day often hidden under the
common human, and recognizable only on closer consideration.'' But what,
then, is the dispute about? Partly about the further development and
carrying through of what has already been attained, and partly about the
deeper understanding and grounding of it. And here the relation is, more
precisely, this: the spirit and direction in which the development is
carried further will depend on the way in which one conceives and grounds
the ethical as a whole. Here the great opposition between the different
views of life makes itself felt.
\opage{2}However much two people may agree about an ethical relation or a
moral commandment, they may disagree just as much in their grounding of it.
And yet in the ethical domain the grounding cannot be separated from the
content. The way in which I ground my action will at the same time be an
expression of the disposition with which it is carried out and of the
motives that call it forth; and in the ethical domain it matters not only
\emph{what} I do, but also, indeed especially, \emph{how} I do it. And
although conscious thought can here, as little as anywhere else, give a
complete picture of what stirs as moving motives in living feeling and
energetic impulse, it is nevertheless a task that inquiring thought cannot
push aside, to give itself an account of the foundation on which ethical
action rests. To make a contribution to the elucidation of this question is
the aim of this little book.

In practical life there is neither time for such deliberations nor direct
use for them. The circumstances of life demand action, often immediate
action. A choice must then be made in accordance with the best insight
there is at hand. It is elements of very different origin that together
make up the practical conviction which determines human action; and however
important it is to undertake an exact investigation and evaluation of all
these elements, such self-examination can in practice be undertaken only to
a very limited extent. Partly it presupposes a clear self-consciousness and
a skill in psychological analysis that are given to few, and partly it is
undertaken not with contemplative calm and disinterestedness, but under the
influence of mighty feelings and passions. Every human being is a practical
ethicist insofar as his striving must aim at bringing his action into
harmony with the highest end and the highest law that he can recognize; but
his endeavor toward the purification and deeper grounding of what is end and
law for his \opage{3}action will soon be checked by limits that lie not only
in his own individuality, but also in the given social and historical
conditions. A young Fiji Islander, for example, cannot conceive otherwise
than that he gives his filial love its most fitting expression by taking his
mother's life while she is still in her vigorous years, so that she shall
not pass into the other world as one decrepit.\footnote[1]{See the
conversation between the young Fiji Islander and the missionary in
\emph{Lubbock}: Origin of civilisation p. 368.} A more far-reaching ethical
criticism would here presuppose a break with the tribe's inherited belief in
immortality. --- A wider horizon we find in those whom one might call the
prophetic ethicists, those personalities who, by virtue of a deep
understanding of the conditions of the people and of the times and of the
fundamental relations of human life as a whole, and by the power over men
that their spiritual superiority gives them, are able to discover and spread
new ethical laws, to create new forms and directions for ethical life. Great
advances in an ethical direction often become possible only through such
prophetic personalities, to whose number the great founders of religions,
lawgivers and reformers belong. But despite the greater breadth and depth of
these minds, there is nevertheless to be traced in them too the limitation
that is inseparable from every striving directed toward particular practical
ends. Their field of vision is conditioned by the historical presuppositions,
by tradition, and by the nature of the people from which they spring and on
which they are to work. Psychological and historical criticism, in them too,
decidedly recedes into the shadow before the imperious demands of action.

If, then, the task we have put forward --- to find the elements of the
ethical and the law of their combination, --- is to be treated at all,
another standpoint must be taken. The theoretical ethicist relates himself
in a purely investigative way to what the practical and the prophetic
\opage{4}ethicist always grasps only in immediate relation to practical
action. If the practical ethicist is the sailor and the prophetic ethicist
the helmsman, then the theoretical ethicist is the astronomer, who from his,
at least relatively, fixed station is able to make observations and
investigations of quite another certainty and exactness than those who
battle with wind and wave. There is no presumption in the position of the
theoretical ethicist toward the other two; if he sees more clearly, it is
only because he moves in the world of ideal conceptions instead of in full-bodied
reality, and if he would set his theoretical insight above the energetic
conviction that they have attained by the practical way, he must remember
the well-known verse: ``To understand a chart is one thing, another to steer
a ship!'' --- % „Et er et Søkort at forstaa, et Andet Skib at føre!“
The theoretical ethicist brings no other presuppositions to the subject he
would investigate than any other inquirer. What matters first and foremost
is to grasp the object one has chosen in its characteristic peculiarity, to
find the most fitting expression for it. As regards the ethical, then, the
first thing will be a description of the psychological state, the relation
and the disposition, that rightly bears its name. The next is to separate
out the elements of which this fundamental form consists and to show the
sources from which they derive; in this way ethical life will at the same
time be seen in its connection with human life and the life of the world as
a whole. Theoretical ethics is thus an empirical science, built on
psychology and history. It proceeds from the observation of action as
actually given and infers from it to general laws. But it cannot remain
standing at the empirical. Every scientific cognition is an ideal
construction, which grasps the elements given in experience while passing
over all disturbing and inessential circumstances. Thus, e.g., Kepler's laws
are ideals to which the actual motions of the planets correspond only
approximately. By abstracting --- \opage{5}i.e., by disregarding all
incidental circumstances --- one at the same time idealizes, i.e., one forms
a picture of the matter as it would be under certain simple presuppositions
and without the disturbing interference of other conditions. In this way
theoretical ethics also forms an ideal type, an expression of human action
as it would be in its perfection. The consistency and purity with which such
a type is developed have a purifying and clarifying effect back upon the
intricate conditions of actual life; and even if it would be doctrinaire
fantasy to apply such a standard directly under certain given conditions, it
is nevertheless of great importance to become aware in which direction the
movement must go. The conception of the actual and the conception of the
ideal thus serve each other for mutual enrichment: the actual is explained
and illuminated and is thereby raised to a higher stage; the ideal receives
closer, more concrete determinations and is thereby brought nearer to
realization.

\begin{center}\rule{6em}{0.4pt}\end{center}
% p. 6 is blank in the print

% ===== PSYCHOLOGICAL REMARKS (pp. 7--18) =====
\chap{Psychological Remarks.}
\opage{7}Our sensations and ideas arise in us as symbols of the nature and
relations of things. When we perceive or imagine something, we look away
from ourselves and become absorbed in the contemplation of the object, and
the more we are able to do this, the clearer and more exact the conception
we gain. But every sensation and idea also acts back upon our
consciousness, intervenes in one way or another in our whole mental state,
produces a certain mood or feeling: --- if our conscious life as a whole is
furthered and heightened, a feeling of pleasure, --- if it is hindered and
weakened, a feeling of pain. These moods can often be wholly subordinate in
relation to the outward-directed side of the sensation or idea; the more
they gain the upper hand, the less energetic and clear the conception of
things becomes. He who is moved by pleasure or pain is no longer an
impartial observer, and conversely, in order to make observations,
\opage{8}pleasure and pain must be pushed back. There is thus an inverse
relation between cognition and feeling; the two cannot both stand at their
highest point at the same time, but the highest degree of the one
presupposes the lowest degree of the other.

To everything that calls forth a heightening of the mental state, and thus
a feeling of pleasure, we ascribe a certain value, which originally shows
itself in our involuntarily seeking to hold fast the idea of it and to move
toward it. The sensation and the idea in themselves contain no ground for
movement; it is only through its relation to our mental life as a whole,
thus through the feeling it awakens, that a cognition calls forth movement
and action, and since an inverse relation obtains between cognition and
feeling, there must also be an inverse relation between cognition and
action.

But this is only one side of the matter. In the preceding remarks we have
taken account only of the \emph{degree of strength}. If, on the other hand,
we ask about that by which the feeling is awakened and toward which the
movement is directed, thus about the \emph{content}, it will turn out that
cognition and feeling, cognition and action, no longer stand in an inverse
but in a direct relation to each other. A feeling of pleasure leads us to
hold fast the idea, or to continue the movement, that calls it forth. But
when certain ideas and movements now turn out to stand in close connection
with one another, the one leading to the other as cause to effect, or there
is \opage{9}great similarity and kinship between them, then the feeling will
gradually be extended to all such ideas as belong together. There then
forms a comprehensive world of objects to which the individual attaches
value, because they directly or indirectly awaken and heighten mental life.
On the content of this world it depends in which direction the individual's
life and action will go. Every higher stage that is reached in the
development of knowledge will therefore necessarily act back upon the life
of feeling and of action. Quietly, but without interruption, feeling moves
up in the footsteps of thought, becomes nuanced, made precise and enriched.
The higher and richer the life of thought, the deeper and fuller,
therefore, the life of feeling can also be. Many feelings that the
civilized European knows are foreign to the savage because of his
little-developed circle of ideas, and the feelings that they have in common
have a different character in each of them.

When experience has attained a certain range and knowledge a certain
development, action takes account not merely of the immediate and present
feelings of pleasure and pain, but also of the possible and more remote
ones. But this already presupposes a stage of development to which several
savage peoples have scarcely raised themselves. They fill themselves with
food and drink when they have abundance, without thinking of the harmful
consequences of gluttony or of the lack of food that will soon succeed the
abundance. They think only of the momentary satisfaction. Only need brings
them to work again, and \opage{10}so they go on in the same cycle. Only when
conscious life reaches a certain development, in strength as well as in
range, can memories and expectations determine the mode of action, even
when it is counteracted by the momentary pleasure or pain; and we have here
the first trace of a rule for action lying beyond what is immediately given
and present. And at the same time we have here the first clear appearance
of a real will, in the energy with which the momentary promptings are
pushed back so that the motives that have proved to be of more essential
importance for the individual's weal and woe can assert themselves. In all
attention there is something of a willing, but this appears most clearly in
the domain of action, as soon as experience has called forth a
consciousness of the different consequences of actions for man.

Attempts have often been made to reduce attention and all manifestations of
the will to mere modifications of sensations and ideas, a kind of
condensation of these. On the contrary, recent physiology and psychology
have led to the result that movement precedes sensation and is a more
primitive expression of our nature than the latter. At the stage of
conscious life it seems, at least, as if movements and actions were guided
by definite motives and purposes. But if we go back to the stage where
there is as yet no question of conscious sensations and ideas, to the child
in the mother's womb, then its voluntary movements cannot be derived from
any purpose, cannot be determined by any \opage{11}idea, but have a purely
spontaneous character. ``It moves,'' as Johannes Müller says, ``its limbs
not in order to attain an external end, but only because it can move
them.''\footnote[1]{\emph{Joh. Müller}: ``Handbuch der Physiologie'' (I. p.
94). \emph{Alex. Bain} has in his work ``The senses and the intellect''
(Book I) more fully developed and grounded the doctrine of spontaneous
movement, which, moreover, one will already find in \emph{Descartes}, in his
``Traité des passions'' (I. 10---11., II. 107).} Only gradually, then, do
the experiences made in this way lead to guiding the following movements in
certain directions, and finally the idea of the possible movement can
precede the actual movement and call it forth. But the very manner in which
the original spontaneous movement takes place will be conditioned by the
individual's innate organization, and from this it is again seen to how
high a degree the direction and character of the individual's action must
depend on the nature that it receives as an inheritance from its ancestors.
Individual action here goes back into the action of the whole race and is
determined by it. For it is not only at the beginning of individual life
that spontaneous movement plays a part; all that we call instinctive action
must be referred to it. These are actions that we perform without knowing
why, and whose real character we only later become conscious of, just as
the young bird sets its vocal organs in motion without knowing what it is
doing, and only afterwards learns to know its own notes. And even when
conscious motives assert themselves, everyone nevertheless knows from his
own \opage{12}experience cases in which neither any single one of these
motives nor all of them together contain the sufficient explanation of the
action; one must then here go back to the deeper instinct, whether this has
formed itself as a second nature in the course of the individual's life or
springs from the original organization.

The proposition that all action springs from pleasure or pain is thus only
partly true; man does not always act on motives, because he does not always
act with consciousness. It is important to hold fast to this, since
otherwise one will arrive at a false conception of human willing and acting
at its primitive stage. For the proposition about pleasure and pain as the
fundamental motives of action has often been used to prove that egoism is
the deepest impulse in human nature. But real egoism presupposes a
reflection, a calculation, and thus a development of the understanding, of
which there is as yet no question at the primitive stages. Instinctive
self-preservation cannot rightly be called egoism. Only at a higher stage
of development can the individual tear itself loose from the connection in
which it originally stands, make itself the center and reduce everything
else to a means. Originally the individual does not make this distinction,
does not distinguish between pleasure and pain and that which calls them
forth, but believes that in the pleasure or the pain itself it comes to
know the nature of things. On this rests precisely the fact that man can be
educated; through the power of changing moods he can be led beyond his
\opage{13}original limitation, which would not be possible if he had really
isolated himself in himself from the beginning.

Just as little as our physical world-system could have developed out of the
original nebulous mass if the repulsive force had been the only one
prevailing, just as little would the world of humanity have come into being
if only one single impulse, the isolating impulse of self-preservation, had
ruled in human nature. In this nature too a plurality of forces is at work.
There lies in it a predisposition to sympathy just as much as to egoism;
there is in our nature ``something of the dove beside the serpent and the
wolf.'' We find sympathetic impulses already within the animal world, in
the parents' love for and self-sacrifice for their young, and in the
cohesion between individuals of the same flock. In man these impulses are
developed and enriched like all others, as spiritual development advances.
Involuntary compassion and fellow-feeling with others cannot be derived
from any calculation of what consequences their happiness or unhappiness
may have for us. Sympathy presupposes only that we ourselves know from
experience what pleasure and pain are, and that we really regard others as
beings of the same kind as ourselves. We involuntarily put ourselves in
their place, make their feelings our own. The extent to which sympathy
asserts itself depends on how far the consciousness of community reaches.
Sympathetic feeling has its original hearth in the family relation, and
among very crude tribes it can easily look, even here, \opage{14}as if it
were overshadowed by the tyrannical impulse of self-preservation, the
individual, in a desperate struggle to subsist, raging blindly against
everything that stands in its way, insofar as superior force does not hold
it in check. But sympathy is in itself just as natural to man as the family
relation is, and just as the human race would perish if % misprint in the print: "Menneslægten" for Menneskeslægten
this relation ceased, the same would happen if all sympathy were lacking.
From the narrower circle of the family it spreads at last to the people,
the tribe, and at last to the whole of humanity.

Both by following the promptings of the impulse of self-preservation and by
following those of sympathy, man is led to regard himself as a member of a
greater whole. Whether his impulses are satisfied or meet resistance, he
receives an impression of definite conditions and fixed connection, of an
order of things under which he must bow. It is, as we have remarked, the
development of knowledge that makes a further development of the feelings
possible. But knowledge does not stand merely as a serving instrument. It
itself makes essential contributions to the character of action. It sets
up general laws and rules, carries through one point of view consistently.
Feeling and passion have a tendency to make exceptions, to act irregularly
and only in particular directions. Developed knowledge makes action
consistent and directs it toward one essential end. The individual senses
in the known truth an authority that stands untouched by his wishes and
impulses; he has in it a fixed point to which he can take \opage{15}refuge
in order to win an impartial conception of his conditions and his conduct.
The personal life needs to be purified by the impersonal truth, raised
above all arbitrariness and contingency.

From what has now been developed we draw some general conclusions, which
will lead us nearer to our proper subject. For man to come to act, there
must in the first place be a sufficiently strong occasion for it. The
strength of the occasion here means the degree to which man's nature is
disposed to action at a given moment, and such a disposition can, as we
have seen, be either an original instinct, an unconscious impulse, or a
feeling of pleasure or pain awakened by certain sensations and ideas. But
in the second place we have seen that experience and knowledge have a great
influence on what awakens pleasure and pain, and that man in his action is
determined by the greater whole of which he is a member, and which he can
embrace with disinterested sympathy. So far as we look merely to the
\emph{strength} of what drives man to action, we seem to have to do only
with himself; but the \emph{content and end} of action point us beyond the
single individuals. The question now is whether in any of the psychological
laws brought out here we have starting points for the conception and
grounding of the ethical.

That the ethical does not rest on the strength of what leads to action
will easily be seen. Strong feelings and instincts are immediate
expressions of the individual's temperament and innate character. It
\opage{16}can be an advantage, in an ethical respect too, to be equipped
with strong natural impulses, but no one will immediately find the essence
of the ethical in them. All will agree that the ethical always presupposes
a law, a rule, an end, which does not immediately coincide with impulse or
instinct. It is with the good as with the true. We do not call an idea or
an opinion true until we have convinced ourselves of its agreement with the
real, according to the standard that we possess in this respect. And
neither do we call an action good except insofar as it accords with the
general rule which in our opinion ought to hold in human life. Here, then,
a higher law or standard is presupposed in both places. We are thus
referred over to the other side of human action, --- its content and end.
Nor will one here immediately find the ethical. There certainly lies, as we
have seen earlier, already in man's immediate nature a capacity to devote
himself and to recognize something that lies beyond his own individual
self. But only when this recognition becomes conscious will it receive the
name of ethical. And this will then be precisely the most general
determination that we can set up of % misprint in the print: "Eestemmelse" for Bestemmelse
ethical action: an action whose motives are ideas and feelings that point
beyond the individual self and attest that the individual sees himself as a
member of a more comprehensive order of things, whose end he makes his own
and whose law he regards as the rule for his life. The Greek word from
which \opage{17}the expression ``ethical'' is formed, like the Latin word
from which the synonymous expression ``moral'' comes, means custom and
usage; in the Danish word ``Sædelighed'' we have the same basic meaning. % "Sædelighed" (morality) kept in Danish: the word is the point
By these expressions we are pointed to a general law and rule by which the
individual's action directs itself, but it does not properly lie in them
that this law is the object of his living feeling and assent. Mere
habituation certainly binds the single person to a greater whole, but can
in itself very well have a purely mechanical origin. Custom is not yet
morality.

Whether this general definition of ethical action, as motivated by
recognition of a higher law and reverence for it, is correct will only
gradually become apparent in the course of our following investigations.
Only in the course of these, likewise, will there be a question of the
nature of the higher law, and of how far we can speak at all of a
universally valid ethical law, as well as of the relation between it and
the various laws by which men at various times and in various places have
felt themselves bound.

The provisional determination of the foundation of the ethical that we have
set up will in what follows be supported first by a more indirect, then by
a more direct argument. For in the first place we shall examine the
justification of the attempts that have been made to derive the ethical
from one or more of the individual starting points brought out in the
foregoing: the pure impulse of \opage{18}self-preservation, sympathy, and
reason. If these elements contain the full foundation of the ethical, our
definition will be incorrect; if, on the other hand, those attempts prove
to be incomplete, it will gain an indirect confirmation; for if the
foundation of the ethical cannot be sought in the purely individual, it
must be sought in the relation in which the individual stands to a higher
power. The second thing we have to do will then be to show how such a
relation to a higher power comes into being, and how the recognition and
reverence develop in which we see the properly ethical motive.
--- The first series of investigations will thus concern individuality, the
second authority.

% ===== INDIVIDUAL STARTING POINTS (pp. 19--54) =====
\chap{Individual Starting Points.}
\opage{19}1. We have already seen in the foregoing that it was only wrongly
that one could designate man's action, as determined by pleasure and pain, as
egoistic. When one tries to derive the ethical from the impulse toward
happiness, from man's original striving for well-being, one need not therefore
deny that there are other feelings and impulses in human nature than the
egoistic ones. It is enough if one starts from what no one can well deny, that
these are from the first the strongest. In order, then, to find a really firm
and reliable ground on which to build the ethical, one turns to the original
fundamental impulse, which in various forms must find its satisfaction in
every human action. Here one is sure of commanding sufficiently weighty
motives; one appeals to individual interest and seeks to show that ethical
action is only a \opage{20}particular way of satisfying this interest. The
ethical is interest well understood.

In antiquity this standpoint was represented by \emph{Aristippus} and
\emph{Epicurus}. The individual's feeling of pleasure and pain is for them the
foundation and principle of all action. \emph{Aristippus} designated pleasure
more precisely as momentary enjoyment. The task for him then became to possess
such a virtuosity that one can at every moment take an enjoying attitude. But
this isolation of the moment from past and future soon proved impossible to
carry through, just as that virtuosity is the affair of only few mortals, who
combine fortunate outward conditions with mental mobility and a sanguine
temperament. Epicurus therefore did not start from momentary pleasure, but saw
the foundation of right action in man's longing for happiness in the sense of
a lasting contentment secured by prudence, which presupposes resignation in
the face of the unavoidable pains. The lesser is sacrificed to the greater,
the shorter to the more lasting. While Aristippus gladly gave himself up to
the present, Epicurus withdrew into himself, in quiet enjoyment of the peace
within him, the unrest of passion being subdued by rational deliberation, and
the fear of the unexpected and of the supernatural removed by insight into the
necessary laws of nature. The Epicurean joins himself to other men only in
order to heighten his enjoyment by sharing it with his friends. Social life
arises from the mutual utility men can render one another; it procures
security and is thus a necessary condition of the wise man's
\opage{21}happiness, since he would otherwise live in a constant unrest and
would be drawn away from himself into strenuous activity directed outward.

The kindred tendencies that have appeared in modern times have, despite the
common starting point, two characteristic peculiarities. The main thing is no
longer to procure for the individual an idyllic refuge where he can enjoy life
in peace, but the satisfaction of a restless craving for progress. Characters
like Hobbes and Bentham have little in common with Aristippus and Epicurus;
they are polemicists and reformers. Connected with this, too, is the fact that
they think far less of the happiness of the single individual than of that of
all individuals. It was natural for Aristippus and Epicurus to withdraw into
the individual self, since they lived in a period in which public interest had
faded and the period of dissolution of a whole culture had begun. But it was
just as natural for men of a kindred tendency in modern times, when cultural
and social life is in vigorous advance and new forces and ideas are stirring,
to adopt a more comprehensive point of view.

The first good, the one pointed out by nature itself, is according to
\emph{Hobbes} self-preservation. All strive to preserve and enjoy life and to
avoid death, especially a painful death. From this first good other goods can
be derived. Thus power, wealth, friendship and wisdom are goods, because they
serve the preservation and enjoyment of life. A highest good, an absolute end,
cannot be conceived; it would be a state in which all opposition and movement
had ceased; but then \opage{22}no enjoyment, not even any perception, would be
possible, since all sensation presupposes alternation and plurality. The
greatest happiness consists in advancing, with as few hindrances as possible,
toward higher and more remote ends.

Since the impulse of self-preservation is equally strong in all, but the one
often stands in the other's way, men live, says Hobbes, in a constant fear of
one another. The weak fear the power of the strong; the strong fear the
cunning of the weak. The only means to security is to forestall the other by
depriving him of the capacity to do harm. Thus arises a war of all against
all. Men are like raging wolves toward one another. In this condition no
culture, no true human life, can develop. Man's capacity to cast his gaze out
into the future only makes him unhappier than the animal. Fear and anxiety
awaken ever anew, just as Prometheus' liver grew again in the night after
being devoured in the day. Now when reason develops, it teaches man that the
state of war is a source of all his misery and therefore ought to be avoided
at any price. Peace is thus a good about whose value all can agree; it accords
with the interest of all and leads them beyond the miseries of the state of
nature. Reason enjoins the natural laws that men must follow if they would
work their way out of the original animal condition. These laws set bounds to
individual freedom, to the immediate influence of the passions, and enjoin
certain virtues \opage{23}as necessary means to a peaceful life. Such virtues
are justice, gratitude, placability, compassion, etc. --- But in order that
the natural laws should be more than mere precepts and doctrines and be able
to come forward with sufficient authority against the unruly impulses, all
must submit to a common sovereign, an absolute authority, before which each of
them entirely gives up his own will. The law of nature, which prescribes the
means for the satisfaction of the impulse of self-preservation, therefore also
prescribes complete submission to the will of the power of the state as the
only means of securing peace.

\emph{Bentham} is not pessimistic like Hobbes; he does not consider it
necessary to hog-tie men so that they shall not tear one another to pieces
like wild beasts; but he nevertheless holds fast to the impulse toward
happiness as the only foundation of the ethical. Against this impulse all talk
of duty and self-sacrifice is only empty words; how should such phrases be
able to overcome the mightiest craving, the one most deeply grounded in our
being? And on the other hand, what kind of virtue would that be which in no
respect whatever contributed to furthering human happiness? And even if there
were such a virtue, by what motives would one then be able to get men to
strive for it? One must reduce all questions of good and evil to questions of
pleasure and pain; only then can one arrive at a sure and definite valuation.
From the simple opposition between pleasure and pain \opage{24}all ethical
virtues, all ethical relations, can be grounded. Give me joy and sorrow,
pleasure and pain, and I will create an ethical world!

Since the fundamental impulse in man is originally directed toward happiness,
it is merely a matter of guiding it rightly; for from shortsightedness and
faulty calculation man often misses what he seeks. The task of ethics, then,
is to teach man to calculate correctly. He must learn to sacrifice the lesser
for the greater. Only such sacrifices can be defended as increase pleasure or
diminish pain, now or in the future. Prudence therefore becomes the first, or
properly the only, virtue. Moderation and self-control stand in close
connection with it. Not every striving for happiness or utility, however,
deserves to be called virtuous; for a striving to bear this name, it must be
bound up with effort and have a resistance to overcome. The effort consists in
holding fast the insight one has gained into what brings pleasure and pain,
against the promptings of the moment. An unethical action rests on a faulty
calculation of personal interest; but the individual has not only his own
interests to take into account. In order to be able rightly to reach his own
goal he must serve the interests of others. A great part of his joys and
enjoyments rests on the will of others, and he can secure their cooperation % misprint in the print: "han kun sikre" (kan wanting); translated as "can"
only by himself working for their wishes. It is personal interest itself, the
strongest of all bonds, that ties the single individual to the whole race. Do
not imagine that men will move the tip of their finger \opage{25}to serve you
when they can have no advantage from it; it has never happened and will never
happen, as long as human nature is as it now is. But men will wish to serve
you when they find their advantage in it, and there are countless occasions on
which they can be of use to you at the same time as they are of use to
themselves. Such occasions the understanding discovers, but the multitude does
not notice them. In these mutual services virtue consists. --- Beside prudence
as the chief virtue is thus set calculating benevolence. But Bentham expressly
insists that love of oneself lies at the basis of love of others, or, as he
expresses it, that the social principle is subordinate to the personal
principle. Ethics, as Bentham conceives it, could be defined as the doctrine
of the detours men are compelled to make in order to satisfy their striving
for happiness. The most important of these detours is conditioned by regard
for others. Bentham lays such great weight on this point that he sets up as a
general principle: the greatest possible happiness for the greatest possible
number of men! This principle does not cancel the command of personal
prudence; for the happiness of each single individual is, after all, just as
much a part of the common sum of happiness as that of all others. The
principle is thus fulfilled no better than when each, for his own part,
procures for himself as much happiness as possible.

Bentham polemicizes vehemently against words like ``duty'' and ``ought,''
because in his opinion they favor \opage{26}ungrounded commands of authority.
When one has declared that something is a duty or ought to happen, one regards
the matter as settled without grounding it further. But there is a little word
that has remedied the harm those words have done: it is the word \emph{why}.
If only one goes on using this, one will, Bentham hopes, put an end to the
belief in authority and to the dominion of mystical inspirations in ethics.
And here he really has a great merit. He has, as energetically as perhaps no
other ethicist, demanded grounds for each single ethical command. He has
further shown the close connection between the belief in authority and
asceticism; both are equally tyrannical toward human nature, and mutilate it
each in its own way. Finally, in his principle of the greatest possible
happiness for all he has set up a standard for practical reforms which has
already proved its great significance. But the central ethical relation he has
no eye for in his theory. If there is anything an ethicist should occupy
himself with investigating, % misprint in the print: "under-|dersøge" for "undersøge"
it must be the very concept of obligation. That Bentham wants to expel this
from ethics has its natural ground in the fact that this concept can find no
place in his own conception; for one need not, after all, be obligated to what
one immediately and involuntarily strives for. On the other hand, Bentham does
not refrain from using both the word ``duty'' and the word ``ought'' when he
wants to enjoin the virtues of prudence and benevolence. And the only sanction
he can then point to as the foundation of obligation becomes the fear of
punishment, either criminal punishment \opage{27}or the condemnation of public
opinion. Here, like Hobbes, he appeals to fear as the final motive.

The general psychological investigations we have undertaken earlier furnish
the necessary elements for an assessment of the attempts now set forth. What
Bentham (we keep to him, since he has given the theory its most finished form)
presupposes is that man's action takes place only according to conscious
motives that can be weighed mechanically against one another. He overlooks the
deeper connection between the single individual and the whole race from which
it springs and in whose midst it lives. According to Bentham's doctrine it is
only personal interest that binds men together. Although he does not share
Hobbes' doctrine of the original contract, by which social life was founded so
that each individual might gain security and peace, he nevertheless still
stands in principle on the same standpoint. Society and the race are for him
only a result, a mechanical aggregate of individuals, not also the maternal
womb from which these spring. But Hobbes himself had already become aware that
there was, after all, one society that could not be explained in this way,
could not have arisen from the individual arbitrary choice of the parties,
namely the society of the family, especially the society between mother and
child. Here not only does an immediate sympathetic instinct assert itself, but
we have here also a relation in which authority can in no way be a consequence
of contract. The child stands immediately under the mother's authority, or, as
Hobbes \opage{28}expresses it in his language, there can be no son in the
state of nature. Man thus in reality never begins entirely from the beginning,
and already in this lies a hint that ethics too, the doctrine of right human
action, cannot be grounded deeply enough by taking a purely individualistic
starting point. Man, as enjoying and suffering, makes everything a means for
himself, directly or indirectly; but the ethical relation first arises when
man feels himself not merely as end but as means, not merely as master but as
servant, --- means and servant for ends that are more comprehensive than his
own individual pleasure and pain. This does not conflict with what Bentham,
with such great right, emphasizes against ascetic ethics, that no sacrifice
has value and can be defended that does not produce more happiness than is % misprint in the print: "forsvares. som" (a point for a comma)
sacrificed. For because the individual's happiness cannot be the foundation of
the ethical relation, it is not thereby said that the ethical stands hostile
to it. It might well be that the highest happiness, the fullest unfolding of
man's capacities and powers, precisely presupposes that he stands in the
service of a higher cause and wholly gives himself up to it without anxiously
heeding his own individual pleasure or pain.

Only when the question is not of the foundation of ethics but of its content
does Bentham's greatness show itself. It is a great philanthropic and
reformatory principle that he has set up in the proposition of the greatest
possible happiness for all. His doctrine is rather a theory of right
legislation than \opage{29}an ethical theory. Legislation must always strive
to bring about a harmony between individual interests. It looks to the result,
not to the motive. Thus Bentham also teaches that it is the action itself and
its consequences, not the motive, that matter. All motives are, after all, in
his view only different forms of one and the same fundamental impulse; men act
only to obtain pleasure and avoid pain; in this, then, all actions agree,
whereas their consequences are very different. Nevertheless, it is not the
consequences in themselves either that, in Bentham's view, make an action
ethical or unethical; he calls only that action virtuous in which a more
comprehensive cognition, and not the momentary prompting, is the determining
ground; virtue presupposes effort. He thus himself goes back to the motive (of
which regard for the consequences must of course always be a part), and
confirms the provisional definition of ethical action that we have set up, in
that he too demands a higher law as determining the action, only that this
law, in his view, springs from the impulse toward happiness itself, is a
direct interpretation of its demands. But this is too narrow a
standpoint.\footnote[1]{In the defense that Stuart Mill has made of
``Utilitarianism,'' or the ethics grounded on the impulse toward happiness, he
has, as I have shown elsewhere (Den engelske Filosophi i vor Tid, p. 51---58),
so transformed this doctrine that he almost no longer defends what its
opponents have attacked. % misprint in the print: a comma for the full stop after "angrebet"
There are two points in particular in which he departs
% [the note runs over to p. 30 here]
from Bentham. He abandons the doctrine of Bentham and of the oldest
utilitarians, that all pleasures are alike in kind and differ only in strength
and duration, and distinguishes between higher and lower enjoyments. But since
the standard for this difference in value cannot be sought in the enjoyment
itself, utilitarianism is properly abandoned with this. He next lays far
greater weight on the sympathetic feelings and sees in them the proper
foundation of the ethical, whereby he gives up the purely individualistic
starting point that was the sole ruling one in Bentham.}

\opage{30}It might seem as if the doctrine of the impulse toward happiness as
the foundation of ethical action were confirmed by \emph{Darwin's} well-known
doctrine of the struggle for existence. Like the plant and the animal, man too
reaches the highest development of his being only through the labor of
preserving and sustaining life. It is a fact that culture and social life are
to a great extent possible only through the interaction of interests. It is
the never-resting striving for a higher and more lasting enjoyment of life
that leads to new labors and discoveries; the savage is held back at a lower
stage of culture by his few and simple needs. But however high a stage of
culture one presupposes, life is nevertheless always a struggle to survive. At
any moment the simple question of being or not being can be posed, and the
impulse of self-preservation can then express itself in all its brutal
nakedness. Thus, e.g., on board a ship that is burning on the open sea, and
where the crew seize the boats while the passengers fight for life and death
over the life belts. The strongest and the most ruthless will here most easily
escape. What only a single terrible event calls forth among
\opage{31}civilized men is standing custom among savages, who, like the
animals, rid themselves of sick and decrepit individuals.

What must first be remarked is that, even if egoism were the only motive in
the struggle for existence, this circumstance would nevertheless not furnish a
sufficient reason for deriving the ethical law from egoism; the question would
then be precisely whether there was no way of getting beyond egoism. But one
must not stare oneself blind at examples of brutality, especially since one
could set against them examples of heroic and disinterested
self-sacrifice.\footnote[1]{On the 23rd of December 1871 the steamship
``Amerika'' burned on the voyage from Buenos Ayres to Montevideo. The captain
had at once jumped overboard with a life belt around his body; the crew saved
themselves in a boat. The passengers had only the choice between death in the
flames or in the waves. People fought with dagger and revolver over the life
belts. A wretch killed a young, helpless wife in order to seize her life belt,
while another passenger, Louis Viala, voluntarily gave up his to a young
woman. The wretch was saved; Viala perished. (Berlingske Tidende 3rd Feb.
1872).} The struggle for existence means essentially the struggle to preserve
the particular form of existence one has attained. The further development
advances, the less it is bare existence that is at stake; men have their life
in certain \emph{forms} which they cannot or will not give up, forms that
stand as the expression of the human, and whose preservation and development
is the end of the struggle. The struggle is for the upholding of \opage{32}all
the true, the beautiful and the good that the human race has produced in the
course of its life. He who sacrifices himself for an ideal end, his art or his
science, or he who lays down his life to save another, fights just as much for
life as the egoist who sacrifices everything and everyone for his own
advantage. The political economists speak of an accustomed mode of living, a
``standard of life,'' as the sum of the demands that the workers make on life,
and at whose level they struggle to maintain themselves. In an ethical respect
too there is a ``standard of life,'' and it is really its preservation that
the struggle is about. Man always fights to uphold what is for him the
highest, and to which he has attached his whole personality. One has no right
to look only to the most elementary and brutal forms of the struggle for life.
To come to know the nature of man, one must not merely study it at a single
stage of development, but follow it through all the different stages, and in
the law that holds for the course of development from the lowest of these
stages to the highest one will have the full expression of human nature. The
cabbage caterpillar and the butterfly are the same animal at different stages,
and only in the law of the development from the one stage to the other does
one have this animal's nature expressed.

2. Over against the attempts to derive the ethical from the impulse toward
happiness we now set the attempts that have been made to derive it from
natural sympathy. In antiquity the \emph{Stoics} brought forward this starting
point, showing how love of \opage{33}mankind has its source in parents' love
for their children, from which it gradually spreads in wider and wider
circles, embraces the whole people, indeed at last the whole human race. Just
as we use our limbs before we have learned for what use we have been given
them, so we have received from the very hand of nature an impulse toward
community with others. In modern times this doctrine has been further
developed by \emph{Shaftesbury}, Hume and \emph{Adam Smith}, and in our own
day by \emph{Auguste Comte}, \emph{Darwin} and \emph{Spencer}. In the two last
named it is brought into connection with the hypothesis of evolution.

When man immediately and instinctively shares in the weal and woe of others
and from his very birth is a member of a society, then the preservation of
society will also stand as an end for the individual. Already among the
animals we find strong social impulses. Darwin cites, after Brehm, the
incident that, when a troop of baboons was put to flight by travelers throwing
stones, some of the strongest males turned back several times to rescue, at
the risk of their lives, the young that could not keep up. Such social
impulses are developed by natural selection. The animals that help one another
gain an advantage in the struggle for life over those that fight in isolation.
Through inheritance and constant exercise the social interests are developed
ever further; and since the survival of the race depends on them, they will as
a rule be the strongest; they work quietly, but without interruption, even if
at particular moments they are pushed back by the egoistic impulse of
self-preservation and of happiness. When the individual, after having thus
\opage{34}succumbed to temptation, has come back to its normal state, it will
notice that a more lasting and more deeply grounded instinct has been hindered
in its activity. Already in the animal, therefore, one can speak of remorse;
yet only man can be called a moral being, because he has the capacity to make
comparisons between his past and future actions or motives and to approve or
disapprove of them. One can often observe in animals a struggle between
instincts, e.g., in the dog between obedience to its master and love for its
puppies, in the swallow between maternal love and the migratory impulse. ``As
long as the mother swallow is feeding her young or brooding over them,
maternal love is probably stronger than the migratory instinct; but the
instinct that is the more persistent wins the victory, and at last, when her
young are for a moment out of sight, she takes flight and deserts them. What
scruples of conscience would not each single bird feel if, when it had come to
the end of its journey and the migratory instinct no longer worked, it were
endowed with so great a liveliness of mind that it could not prevent there
constantly passing before its soul an image of its young, dying in the wintry
north of hunger and cold!''\footnote[1]{\emph{Darwin}: Om Menneskets
Oprindelse. Danish trans. I. p. 86.} Man's intellectual capacities are now
precisely so developed that earlier impressions and images constantly pass
before his soul, so that he cannot avoid making \opage{35}comparisons; only
he, therefore, can feel real remorse; it awakens in him a sharpened feeling of
duty, i.e., of obedience to the most deeply lying and constantly working
instinct in his nature; that there is something man \emph{ought} to do means
that he has a consciousness of the existence of a persistent instinct which
serves him as a guide, but is, to be sure, liable not to be obeyed.

Just as, according to Bentham's doctrine, not every satisfaction of the
impulse toward happiness could be called ethical, but only that in which man
followed the prudent calculation of the understanding despite the momentary
prompting, so the adherents of the morality of sympathy cannot regard every
action from sympathetic impulse as ethical either. Tacitus relates of the
ancient Germans that, although they took pleasure in gifts, they thought
neither that they owed gratitude when they had received them nor that they had
any claim when they had given them. And likewise it is related of some Indian
tribes that they acted from affection just as they acted from desire, without
regard to the consequences. When they had shown another a kindness, they had
satisfied an impulse; the matter was finished and vanished from memory; but
conversely they felt no regret either at having satisfied their selfish
passions; both in their sympathy and in their egoism they acted as true
children of the moment. Here again it is shown that the ethical presupposes a
model, a law, which lies beyond the mood of the moment and points toward a
more comprehensive and fixed form \opage{36}for human life. Therefore Darwin
and his predecessors emphasize the significance of memory and thought. With
them it stands not in the service of the impulse toward happiness, but in that
of sympathy, and holds fast its demands against the claims of the impulse of
self-preservation and of egoism. The personal principle is here subordinate to
the social. Through the purifying and illuminating power of cognition the
sympathies can gradually be freed from their limits, go beyond the narrower
circles of the family and the nation, and become a disinterested love for all
living beings.

Bentham and the utilitarians too could point to a great social connection in
which the individual stood in an ethical respect. All that belongs to the
sustaining of human life and to the satisfaction of its needs is not brought
about anew within each single generation, but each new generation works on the
basis of fixed institutions, forms and traditions which it inherits from
earlier generations. Culture rests on tradition. It is not only by his reason,
his capacity to combine and foresee, that man raises himself above the animal;
progress would not be great if each had to begin from the beginning. Language
and tradition make possible the accumulation and diffusion of inventions and
discoveries. The system of interests to which utilitarianism refers thus
embraces not only all men living at the same time, but also the countless
series of earlier generations. As beings of culture we are bound in solidarity
with our predecessors. But according to the morality of utility, culture is
only an outward connection
% continues on p. 37
% continues from p. 36
\opage{37}connection, into which the individual lives himself because he finds
his own interests satisfied by working toward its consolidation and
development. Bentham and his school do not look to the influence that culture,
according to the general laws of life, has and must have on human nature
itself. Human nature is developed and transformed through the work of cultural
history. With the changed conditions of life the original impulses and
predispositions will also gradually change. Not only through external
tradition does the individual come to share in what the race has acquired; in
his innermost nature he carries the after-effects of the exertions of the
past, according to the general physiological law that function acts back upon
organization. The original conditions of men had to make them egoists; the
milder forms that the struggle for life gradually assumes must also make human
nature milder. The great spokesmen of the hypothesis of evolution, Darwin and
Spencer, therefore come forward against Bentham, Mill and
Bain.\footnote[1]{\emph{Herbert Spencer} wrote in the year 1846 an essay on
sympathy in the same direction as Adam Smith's ``Theory of moral sentiments,''
which he did not know at the time. In 1850 his ``Social Statics'' appeared,
where the same subject is treated. He also comes back to it at the end of his
Psychology, but his ``Principles of morality'' are only to form the conclusion
of his great work.} Whereas the utilitarians presuppose the impulse toward
happiness as equally strong at all times, the starting point of the morality
of sympathy is not a constant magnitude, but varies continually. The
sympathetic \opage{38}impulse has its deepest foundation, as has been
indicated several times, in the maternal instinct. Auguste Comte therefore
finds the first trace of moral and social life at the point within the animal
world where hermaphroditism has ceased and the sexes are separated. From the
family relation it then spreads in wider and wider circles. The purely
empirical conception cannot sufficiently ground this quiet growth of sympathy.
Hartley, to be sure, and after him in our day Stuart Mill, have rightly
emphasized that by virtue of the laws of the association of ideas man will
come to love for its own sake that which he was at first interested in only
because it satisfied his self-love, just as we involuntarily ascribe to money
a significance and a value that properly belong only to what we can obtain for
the money. He who is constantly compelled to work together with others in
order to attain his own ends will gradually forget his original motive and
feel a disinterested interest in the welfare of others and in the common goal.
The art of education consists precisely in attaching the individual's pleasure
and pain to the right objects, binding his interest to them in such a way that
he can at last be wholly absorbed in them. But the hypothesis of evolution
teaches us that it is not education and individual experience that are to do
everything here. The sympathetic impulse is not merely a product of the
associations of ideas. It grows through the generations, since the ever more
frequent occasion to express itself, and the ever more comprehensive circle
within which it moves, must necessarily also strengthen the predisposition in
individuals to let \opage{39}themselves be governed by it. Quietly and
imperceptibly, then, according to this hypothesis, a metamorphosis of human
nature takes place. Every enthusiastic mood, every noble and philanthropic
feeling adds an atom to the building of the perfect humanity. Nothing is lost,
in spiritual nature any more than in bodily nature.

The ethical law, as it takes shape from the standpoint of the morality of
sympathy, would have as its content the development of perfect, comprehensive
love of humanity. But the question cannot be kept away, wherein the ethical
element in sympathy really lies. As an immediate impulse it can attach itself
just as well to the worthless or even the corrupting as to the good and right.
Adam Smith himself remarks (what would perhaps no longer fit our day) that the
great mass is a disinterested admirer of everything unusual, especially great
power and wealth. Love is, after all, on the whole blind. In order to make
others happy it must have a model, an ideal of happiness, of perfected human
life, toward which it can then seek to lead others. But this it can have only
when thought is more than the deputy of sympathy and can really come forward
as an ordering and guiding principle. Love must subordinate itself to justice
(which according to Leibnitz is caritas sapientis, \emph{the wise man's}
love), as the guiding, distributing and harmonizing principle in human
society.

Adam Smith has made an interesting attempt to show how the influence of
sympathy leads to \opage{40}the formation of an ethical standard. Since we
immediately feel pleasure or pain at the actions of others, even if they do
not concern us, we also come to express our approval or disapproval of them.
We suffer with the victim of a cruel or faithless action and feel antipathy
toward the one who commits such an action. On this sympathy and antipathy
rests our approval and disapproval, which thus arises quite instinctively, as
we involuntarily put ourselves in the place of the others, of the agent and of
the sufferer. We play the part of spectators who feel with the persons
appearing in the drama. Our first moral judgments we thus pass on others. Yet
there comes a point in time when we see that we ourselves are acting persons
in the drama of life just as much as the others, for which reason we must
assume that others judge us no less than we judge them. And sympathy itself,
which seeks to even out the difference between our own state of feeling and
that of others, since every disharmony here is bound up with a painful
feeling, will lead us to examine our nature and our actions in order to see
with what eyes we can expect others to regard us. We thus again put ourselves
in the place of others, no longer merely in order to rejoice or suffer with
them, but in order to view and evaluate ourselves from their standpoint.
Everyone thus divides himself into two persons, a judging spectator and an
acting being. The natural standpoint is abandoned; the individual looks away
from himself precisely in order to come to know himself. It is a fact of
experience that he who in the presence of others falls into passionate
\opage{41}agitation involuntarily tones it down when he feels that the others,
as impartial spectators, do not share it. On this rests the educating element
in living together with others. Everyone feels himself the object of a
constant and attentive criticism and involuntarily arranges his action so that
he can stand before its tribunal. The child, the youth and the man thus learn
in the various circumstances of life to look away from the individual
standpoint. At last the inner spectator and judge in our breast no longer
needs to be strengthened by the presence of outer spectators; we have gone
through the great school of self-command and have our rule, our model, within
ourselves. We now also look away from the finite and empirical standard that
is given in the opinions and feelings of our surroundings, and turn our gaze
toward an ideal of perfection that has gradually formed itself in our mind.
The outer spectator thus first becomes an inner spectator, and this in turn an
ideal model.

When development has reached this point, the original starting point,
according to Adam Smith, passes wholly into the shade. Both the sympathetic
and the egoistic impulses step aside for a higher feeling. ``Since we are
always affected far more strongly by what concerns ourselves than by what
concerns others, what is it, then, that on every occasion brings the noble
man, and sometimes even the base character, to sacrifice their own interests
to the greater interests of others? It is not the mild power of love of
humanity, not the feeble spark of benevolence that \opage{42}nature has
kindled in the heart of man, that is thus able to counteract the strongest
impulses of self-love. It is a stronger power, a more forceful motive, that
acts on such occasions. It is reason, thought, conscience, the inner man who
dwells in our breast, the great judge and arbiter of all our conduct. It is
this that, whenever we are about to act in such a way that we hinder the
happiness of others, calls out to us, with a voice that can drown out our most
clamorous passions, that we are only one of many, in no respect better than
anyone else, and that when we so brazenly and blindly prefer ourselves to
others, we rightly become the object of indignation, abhorrence and
execration. From it we learn our real littleness and the littleness of
everything that concerns us, and only through the eye of this impartial
spectator can the natural illusions of self-love be corrected. It shows us the
beauty of magnanimity and the ugliness of injustice, the rightness of giving
up our own highest interests for the still higher interests of others, and the
ugliness of inflicting on others the least injustice in order to obtain the
greatest advantage for ourselves. It is not love of our neighbor or of the
human race that on so many occasions spurs us to practice these divine
virtues. It is a stronger love, a mightier feeling, that generally expresses
itself on such occasions: the love of what is noble and honorable, of the
greatness, dignity and loftiness of our character.''\footnote[1]{Theory of
moral sentiments. Vol. 1. part. 3. chap. 3.}

\opage{43}By this account Adam Smith has acknowledged that the sympathetic
impulse in itself is only one element contributing to the formation of the
ethical law, but does not contain the full explanation of it, and in
particular does not fully express the feeling, the disposition, that the
ethical model awakens. One will also easily see that sympathy is too weak and
too narrow an expression for the relation to the educating powers in the
history of man and of the race. Sympathy presupposes a relation between
equals, but the educating power stands, at least originally, as a superior. To
that extent Hobbes and Bentham are more in the right when they point to fear
and hope as the feelings that stir in the face of the authorities. However,
this consideration we must hold back until we come to the special
investigation of the concept of authority. On the other hand, Smith's doctrine
of ``the impartial spectator'' invites us to investigate the significance of
cognition, of reason, in an ethical respect. It plays, as we have seen, a
great part both in the morality of utility and in the morality of sympathy,
but in both as a subordinate principle. In order to grasp its significance
better, we now turn to the attempts that have been made to make it itself the
foundation of the ethical.

3. Both the morality of utility and the morality of sympathy concede that the
immediate impulses cannot be called ethical, but that there is first a
question of ethical action when a consciousness has formed of a law and rule
that points beyond the moment's \opage{44}immediate promptings. Such a
consciousness presupposes cognition, reason, thinking. Therefore it is only in
the human world that the ethical is at home; the degree of development of
conscious life that it presupposes is not yet found in the animal world, and
has surely also been attained by men only after the lapse of long ages. It was
long before human life began to be guided by thoughts, and we have not yet got
beyond the beginning. --- When we here set thought or reason, as an ethical
starting point, alongside the impulse of self-preservation and sympathy, we no
longer conceive it as servant, but as master, and proceed from the fact that
conscious cognition can, under certain conditions, determine man's action. The
question is then whether this is an exhaustive explanation of the ethical.

But here we are stopped at once by a misgiving. According to our general
psychological considerations, cognition in itself cannot set the will in
motion. Thought becomes a motive only by calling forth feeling and impulse.
Because I have quite distinctly and clearly seen something, it does not follow
that I act accordingly; I do so only when the insight awakens living feeling
and desire in me. There will certainly be a tendency to bring about harmony
among the various elements of conscious life, but this tendency can in itself
be just as instinctive as the impulse toward happiness and sympathy. % misprint in the print: "Lyksalighedssdriften" (ss)
If we imagine a clearly and many-sidedly developed reason, in accordance with
which the will immediately \opage{45}acts, such a phenomenon will nevertheless
not be referred to the ethical. The ethical turns out here again to presuppose
a distance, because it sets a task that is to be solved, a distance that is
filled out by the idea of a law whose validity man recognizes. It will
therefore also turn out that when reason is set up as an ethical principle, it
is really not reason in and for itself that is meant, but reason as the
expression of man's true nature, so that the demand is that the nobility and
dignity of human nature be given its due.

The general theme of this standpoint has been stated by \emph{Pascal} in a
famous utterance. ``Man is the weakest reed in nature, but he is a thinking
reed; the universe need not arm itself to crush him; a little smoke, a drop of
water is enough to make an end of him. But even if the universe crushes him,
he is still nobler than his murderer, because he knows that he dies, while the
universe itself does not know the power it has over him. All our dignity,
then, consists in thought. In thought we must place our pride, not in time or
in space, which we cannot fill. Let us then labor to think rightly; that is
the principle of morality.'' What Pascal expresses here is the enthusiasm at
belonging to the world of spirit. Thought gives man a dignity that unconscious
beings do not have.

Already in \emph{Socrates} we find this principle. He taught that virtue is
knowledge, because only that action is right which springs from real insight.
Above \opage{46}all, everything depends on knowing oneself; for only he who
knows his own capacities, powers and impulses can conduct himself in the right
way in life. The more detailed grounding that Socrates gives of this goes in
the direction of the doctrine of utility, and Stuart Mill has accordingly been
able to claim him for it. Socrates, namely, impresses the importance of right
self-knowledge by pointing to the unfortunate consequences of ignorance and
delusion. But even through the most prosaic rendering of Socrates' teaching it
nevertheless emerges that he really laid the weight on another point. He
maintained that ``nothing was stronger than insight,'' that it was ``something
beautiful, which is able to rule over man, strong enough to help him through
life.'' He who does not let himself be guided by cognition, but by blind
desires and impulses, is slavish or bestial. A slave is he who has no share in
the beautiful and good. ``Does it not seem shameful to you,'' he says to
Aristippus, ``that a human being should behave like the most irrational
animals, which are driven by sensual desire for the bait and blindly walk into
the trap?'' In his defense speech he likewise points to thought and its power
as the expression of man's true nature, the cultivation of which must lie
closer to his heart than the care for the bodily and for external goods. ``My
friend,'' so he addresses his fellow countryman, ``you who are an Athenian, a
citizen of a great city, the most famous of all for culture and strength, are
you not ashamed to care only for amassing wealth and acquiring honor, while
you do not concern \opage{47}yourself about truth and wisdom, about your soul
and its education!''

Already here, in the very founder of ethics, we thus see that when reason is
set up as an ethical principle, what is really meant by it is the inner
freedom and harmony in man's nature, which is possible only when his action is
in accordance with his highest cognition. Man is then not drawn to different
sides by conflicting desires and impulses; nor is he led blindly by the
impressions from the surrounding world; unity, coherence and independence come
into his nature. He becomes a little world for himself, independent in
relation to the external world, and encloses within himself a harmonious
fullness.

What Socrates had indicated in all simplicity remained the main thought in
Greek ethics. Plato, Aristotle and the Stoics each in his own way conceived
man's ethical action as consisting in his acting according to the highest law
of his nature, with rational insight ruling. They find freedom in accordance
with right insight, and in this freedom true humanity. Precisely this
immediate conviction of the dominion of thought, which brought Socrates to
reject with indignation the assumption that ``anything could be stronger than
insight,'' constituted the one-sidedness of Greek ethics. It overlooked the
manifold conditionedness and limitedness of human spiritual life, and all the
oppositions and contradictions, struggles and defeats, that follow from it.
For it, the will fused partly \opage{48}with cognition, partly with immediate
impulse.

In modern times Kant again brought forward the idea of inner freedom, pointing
to the law of reason as that which unconditionally obligates man. He showed
clearly that the law the ethical presupposes cannot be one immediately given
in nature, but must be owed to another source than that from which man's
involuntary actions spring. The ethical law, according to Kant, cannot be
grounded --- neither psychologically, physically nor theologically grounded.
It points beyond everything given, beyond all experience and reality, in that
it prescribes what ought to happen. Only in reason do we have such a power
pointing beyond all experience. Reason's unconditional and ungroundable
command stands as a fact for consciousness, an immediate demand to which man
must bow. It would lose the character of an unconditional command (categorical
imperative) if it were supported in any other way than by reason's immediate
demand.

The progress from Socrates to Kant is easily seen. Socrates is still able to
ground the value of insight only by pointing to its external utility; this
grounding stands, with him, side by side with the appeal to man's dignity as a
free rational being. Moreover, the thought of an opposition, a duality within
man's nature, is foreign to Socrates. With right insight everything is given;
thought does not appear as law, because an immediate harmony is presupposed
between it and the other \opage{49}mental elements. Kant, on the other hand,
raises reason as ethical principle above everything immediately given, and % the print doubles "og" over the line end
rejects every empirical grounding. And he determines reason as legislative,
whereby a distance is presupposed between guiding and obeying elements in
man's nature.

One would not be wholly wrong in deriving this opposition between Socrates and
Kant from the great religious movement that lies between these two greatest
figures in the history of humane ethics. Christianity set up an ideal fetched
from a supernatural world, exalted above everything that human thought should
be able to reach by its own powers. Do we not seem to hear an echo of this
when, according to Kant's doctrine, the law of reason disdains all grounding
fetched from reality? Christianity declared human nature sinful and corrupt.
Are we not reminded of this by Kant's presupposition of the distance between
the command of reason and man's actual action?

Yet the complete explanation will not thereby have been given. However
profoundly that religious movement has acted upon European spiritual life, one
must nevertheless not overlook that it is itself only one of the forms of the
general opposition between antiquity and modern times. For the Greek, life
stood as immediately given and complete. History did not yet tell him
distinctly and urgently of great conflicts between different stages of human
development; he did not have the conception of life as a struggle for
existence and \opage{50}advancement that modern times have won through their
experiences, their action and research. When life, instead of standing before
the mind as a work of art in a limited frame, turns out to be movement and
development, the mind must turn with a quite different energy forward toward
what is not yet but is to be. The thought of a task now becomes quite
differently alive. Thought hastens ahead, anticipates the new stage of life in
an ideal way, and passes judgment on the earlier stage. The ideal view that
thus forms itself does not have its full explanation in the immediately given
and present. We are here directed to the original energy of conscious life, by
which even the simplest sense-sensation becomes more than a passive receiving,
and by which every freer relation to the content of sensation and experience,
and every working-up of it, becomes possible. It is this energy that makes it
possible to form complete and finished images, ideals, out of the always
incomplete elements in the actual world. Here is the source of all poetry and
art, of all higher ideas. If one will call this capacity in man to form ideal
ideas reason, then it is entirely correct that only through reason, as that
which goes beyond the given, does the ethical law become possible, and that
this law always stands exalted above actual action as a model which the latter
strives to approach.

But when we take reason and law of reason in this sense, we cannot, like Kant,
stop \opage{51}at reason's command as a mere fact; we must investigate how the
idea of such a law has developed psychologically and historically. This task
we cannot evade without violating the simplest rules of scientific inquiry.
The ``facts'' in the spiritual and ethical domain must be subjected to
examination and grounding just as much as all other facts. We must here invoke
the critical Kant against the dogmatic Kant; and he himself comes to meet us
here by his account of how the law of reason determines man's action. Only
through feeling, as we know, can reason awaken the will. The law of reason
stands, by its unconditional authority, as a power exalted above the
individual. The sensuous impulses are pushed back, the immediate feeling of
life and self-confidence are checked and humbled by the thought of this
unconditional law. Man feels before the law within him, just as before the
clear starry heavens, his littleness, at the same time as he is filled with
admiration and reverence; for the ethical law acts repellingly on the sensuous
and egoistic impulse only in order to exert so much the greater attraction on
man's proper self, which feels itself one with the law. In the law we
therefore find an expression of our true will. But whatever one may think of
reason, --- a feeling such as admiration and reverence must, after all, have
its natural history, and we have not rightly understood the ethical until we
have explored it.

Since Kant regards the ethical law of reason as \opage{52}raised above all
experience, everything given and present, he cannot give it any definite
content. But it lies in the very concept of a law that it is universally
valid, does not adjust and accommodate itself to the single individual, but is
common to him and all others. The general content of the ethical law is then:
Everyone shall act in such a way that the principle he follows in his action
can serve as a general rule for all. Man shall regard each of his actions as a
link in a higher natural connection, which is to be realized through his
striving. --- This principle Kant himself regards as purely formal and
independent of experience, and his critics as a rule grant him this, only that
they consider it a fault, he a virtue. But whence, one must ask, does pure
reason get the idea of a universal legislation, if not through experience? In
actual life we see the single human being as a link in society, in the race;
we see definite, factual laws holding for the life of men together, and we
then form, by abstraction and idealization, the idea of a perfect society
where each individual fits himself in with full regard to all others, so that
one and the same law holds for all. In the ethical principle Kant sets up we
can therefore see nothing purely a priori, but a historical idea brought back
to ideal simplicity. As we have seen earlier (above S. 4---5), it lies in the
nature of our cognition to operate with ideal concepts, but these are applied
in ethics in a different way than in natural science. In the domain of nature
\opage{53}we regard the ideal concepts as approximations to the actual
connection, but in the ethical domain the actual as an approximation to the
imagined ideal. Experience always lies at the base; only from it can the
elements for the ideal constructions be fetched.

When we see the matter in this way, we also have points of attachment for the
ideas of the doctrine of utility and of the doctrine of sympathy. Culture,
regarded as resting on a system of interests that mutually support one
another, gets its place in the idea of the ideal society as its real
foundation. Without having its firm ground in the primitive demands of life,
the ideal society hovers in the air. Against every ascetic way of looking at
things it is also of importance to assert the individual's right to
self-satisfaction. What utilitarianism overlooks is that this right, ethically
seen, must have its ground precisely in the individual's dignity as a member
of the realm of human personalities, just as every juridical right, after all,
points back to society and its law as its foundation. The doctrine of sympathy
has had a deeper understanding of the individual's inner relation to the race;
and it rightly emphasizes benevolence and love of humanity as the living bond
that holds society together and makes it more than a mechanically joined heap
of atoms; but since our recognition of the law is a main element in ethical
action, we cannot stop at instinctive feelings. Sympathy as well as individual
interest must be subordinated to the highest understanding we are able to gain
of \opage{54}the fundamental law, end and tasks of human life. The realization
of the kingdom of humanity must stand as the end and fundamental law of our
action.

But here it has once again turned out that we cannot understand the ethical as
long as we stop at the consideration of the single individual. Here too we are
directed back to the great connection in which the individual arises and
develops, and to which he owes what he is and has. Whether we proceed from
reason, from sympathy or from the impulse toward happiness, we are of
necessity led out beyond the domain of individuality. From the maternal womb
of society the single individualities spring, and it is, on the other hand, as
representatives of and co-workers on the ideal society of harmonious
personalities that they reach the highest development. We must therefore begin
a new investigation of our question, such that we no longer proceed from the
single individual but from the higher power that from the outset encloses all
individuals.

% ===== AUTHORITY (pp. 55--86) =====
\chap{Authority.}
\opage{55}The more intricate and complex a subject is, the more necessary it
is to employ the descriptive method. As long as we kept to the single
individual and the general psychological laws that determine its action, the
relations were easy to survey and to analyze; but indispensable as the
individual starting points are, it has nevertheless turned out that the
analysis brought to light essential elements which find their explanation only
when we take a social starting point. The single individual we could in the
main regard as a fixed type; society, on the other hand, and in general the
powers under whose shelter the individual comes into being and develops,
belong to the rich and manifold world of history, where everything is subject
to movement and development. Ethics here points back to history. It has not
been the least cause of the imperfection of ethics that the great significance
of history as its foundation has not been heeded. In particular, the
eighteenth \opage{56}century, whose sons the greatest modern ethicists (Ad.
Smith, Bentham and Kant) were, could not, because of its whole abstract and
individualistic tendency, duly appreciate this significance. It is only our
century, whose open historical sense has also led to the conviction of the
importance of the historical foundation for ethics. Schleiermacher and Hegel,
Comte and Spencer show themselves, each in his own way, imbued with this
conviction.

It has turned out in the preceding that one cannot determine the nature of
ethical action except by referring to the idea of a law; but the concept of
law we take originally from human society. It goes in the ethical domain as in
the spiritual domain in general: we use expressions and ideas that are taken
from the outer world in order to make clear to ourselves the nature and
coherence of the inner world. Obligation, too, is properly a juridical
expression. Both law and obligation presuppose an authority that not merely
commands but upholds the carrying out of the command. Without power no law and
no obligation. But power alone would call forth only fear. It is only the last
resort (ultima ratio) for upholding the law. The authority on which the law
rests appeals to other elements in man's nature than fear. The deterrent plays
its greatest part at the lowest stages of development, where the law is
essentially a fence against danger and ruin. According to the old myth, the
first law that went out to man was the prohibition against eating of the fruit
of a certain \opage{57}tree, since it would bring death. The child first
learns to know the law when it is to be brought up to cleanliness, and here
the averting and the deterring play the greatest part. The law acquires a more
positive significance when the authority stands not merely as a threatening
power, but through its whole superiority comes to stand as a guiding model for
man. When man sees, in the authority to which he is subject, capacities and
powers that far surpass those he feels in himself, he feels sympathy with it,
and the impulse to imitate is awakened. Already among the animals the swiftest
and strongest becomes the leader of the herd. As soon as intelligence is
tolerably developed, admiration is awakened for everything great, new and
striking; and with it reverence and respect combine in turn, the more man,
through his increased experiences, feels his littleness and powerlessness, the
more, that is, his self-esteem is lowered by the contemplation of the model's
superiority. The benefits he receives from the higher power awaken gratitude,
and at last he ties his whole weal and woe to it with free trust and devotion.
Over against mere power man will always be able to think of resistance, until
he is utterly overwhelmed and paralyzed by fear; but when the feeling of
dependence is purified by being combined with the feelings of sympathy, of the
impulse to imitate, of admiration, of reverence and of gratitude, and thereby
develops into recognition and self-surrender, we have the essence of
obligation \opage{58}in the ethical sense. The center of gravity begins to
shift from the outer to the inner.

This relation to an obligating authority can be traced through the whole
course of mankind's development under different forms. An absolute state of
nature without authorities does not exist. Already Hobbes remarks, as we have
cited earlier, that no one is born in the state of nature. Wholly without
authority would be only such beings as grew up out of the earth, as in the
myths of the ancients, or, according to the Homeric expression, were the
offspring of a stone or a tree; and only they could be complete
individualists. The authority of the family is thus the most primitive. At the
lowest stages at which human beings are found living (according to various
recent travelers as well as to Herodotus's account of several savage nations),
the family admittedly does not yet exist as a closed group; and the mother
takes care of the child only as long as it is wholly helpless; but already
through the protection and the benefits that here fall to the child's share,
and without which it could not live, the ground is laid for a relation of
piety which can become the germ of further development. When the roving life
is replaced by more fixed and regular conditions, family life comes to play a
greater part. The father's authority now replaces the mother's. At this stage
marriage is grounded on power, not on love; wife, children and slaves must bow
in equal degree to the will of the head of the
family.\footnote[1]{\emph{Lubbock}: Origin of civilisation.} \opage{59}The
protection the weak enjoy, and the constant living together, in part also
working together, call forth other feelings too besides fear, but this feeling
nevertheless always lies at the bottom. The unconditional authority of the
head of the family within his circle maintained itself even after the family
had become part of a more comprehensive society. In ancient Rome the power of
the father of the family and of the husband was subject to no restrictions
whatever; he was answerable to no one on earth. How profound an influence the
relation to this authority must have exercised on the ways of thinking and
acting of the growing generations will easily be seen; and however changed
these relations may be in modern societies, most people still have within the
family the first objects of fear and love, admiration and gratitude, respect
and reverence. It is under the influence of these living forces that the
circle of ideas is formed and the first rules for action, the first moral
laws, arise. The individual here has his practical world-view grounded, and
the validity of this world-view he feels as one with the very authority from
which it stems. Violation or contradiction of it is felt so strongly because
it strikes the very innermost and deepest feeling in the individual.

The bond of blood is the first that unites human beings and thereby creates an
authority. For long ages it has been the only one, or in any case the
mightiest; but the family points beyond itself, in that the head of one family
stands independent \opage{60}over against that of the other. Through the
fusion of the individual families and clans a civil society arises. An
absolute boundary between the two kinds of society can, historically regarded,
be difficult to draw, since the head of the family earlier had, after all,
much of the authority that the head of state later acquired. In the Roman
state one can trace the original union of the individual clans. ``The Roman
state,'' says Mommsen,\footnote[1]{Römische Geschichte. I. S. 59.} ``rests on
the Roman house, with respect to its elements as well as with respect to its
form. The people arose as the old clans in one way or another joined together;
the Roman territory arose from the united districts of these clans; a Roman
citizen was he who belonged to one of these clans.'' The bond of blood here
plainly lies at the bottom. Such a primitive society can grow only by
adoption; even the relation of friendship first acquires its significance
among many savage and barbarian peoples when the friends have mingled blood
and thus, at least symbolically, brought about a natural bond.\footnote[2]{Cf.
\emph{Maine}: Early history of institutions. S. 228.} Anyone who was not tied
to the tribe by this bond was an enemy or a slave. Since the heads of the
individual clans stand as equals, there will be no permanent head for the
whole horde. Just as the cohesion in general is still so weak that the horde
scatters and each looks after himself when distress is at hand, so too one has
a leader only \opage{61}as long as one needs him. Yet the old and the wise,
the strong and the brave will stand in greater esteem than the others. And the
more the struggle with other tribes makes cohesion necessary, the more will
authority also be consolidated. War therefore seems, as Bagehot in particular
has emphasized, to have worked as a binding means within the tribe. The
leaders in battle were regarded as the protectors of all. Under very simple
conditions of culture no great outward distance can arise between those who
command and those who obey; therefore we first get a properly political % the print doubles "og" over the line end (p. 61)
authority when inequality in the distribution of property removes the rulers
from those who obey them, and when at the same time society becomes large
enough for the ruler to be withdrawn from the constant gaze of individuals.
The ruler now becomes an object of reverence; envy is excluded by the very
greatness of the distance. His will is unconditional law.\footnote[1]{Cf.
\emph{Adam Ferguson}: An essay on the history of the civil society. Part II.
Sect. 3, and \emph{Herodotus's} story of Deioces (Herod. I. 96).} Although
habit and instinct, convention and custom form a kind of fundamental law which
even the despot cannot or will not break, it is nevertheless a sign of
progress when a personal, conscious will asserts itself; only thereby does law
in the sense of a general, impartial and inflexible rule become
possible.\footnote[2]{\emph{Maine}: Early history. S. 392.} And this is a
progress not merely in a purely cultural-historical but also in an ethical
respect, insofar as blind habit is replaced by conscious submission.

\opage{62}Power, the protecting and the commanding power, is the first
foundation of political authority, just as of authority within the family.
Supported on this foundation, the higher and freer forms of social life can
develop. One often imagines the life of savages as free and unbound; but on
the contrary it is led under a heavy pressure of authority. ``They are
governed,'' says Lubbock, ``by rules and conventions that form one of the
cruelest tyrannies that has perhaps ever existed on the face of the earth. The
will of the weak, their property, their life are subject to the dominion of
the stronger. The tendency of the whole system is to give everything to the
strong and the old to the detriment of the young and the weak, but especially
of the women.''\footnote[1]{Origin of civilisation, S. 434.} When this heavy
pressure is taken away, the most complete anarchy, the wildest licentiousness,
often sets in, as among some Negro tribes, where everything is in dissolution
at the death of the chieftain. It is a severe school that the human race must
go through at this stage, but it has been necessary. And even the most
despotic state authority must nevertheless have awakened other feelings than
fear; for it appears not merely as commanding but also as protecting, affords
a shelter under which life can go its way. The privileges that the ruler
assumed were in reality bought with services. A great example of this we have
in the institutions of the Middle Ages, the foundation of ``l'ancien régime''.
The king, the nobility and the clergy have in rude and savage times, through
\opage{63}their physical and spiritual authority, made human life and human
development possible.\footnote[1]{Cf. H. \emph{Taine}: L'ancien régime. Livre
I. Chap. 1.: Origine des priviléges.} In the trust and the gratitude that the
protecting power awakens lies the germ of the ethical character of authority,
which is originally pushed back as long as fear and instinctive custom reign.

We have brought forward some basic features of the oldest history of
authority, as it is accessible to us especially through the comparative % misprint in the print: "sammenlig-|lignende" (a syllable set twice)
study of primitive forms of society in recent years. Here the relations are simplest
and easiest to survey; at later stages of development so many different
elements play a part that it is attended with difficulty to consider each of
them by itself. Just as the study of the fetal state affords the physiologist
valuable contributions to the understanding of the completed organism and its
conditions of life, so the study of those simple and primitive forms of life
is also of importance for the psychologist, the ethicist and the
sociologist.\footnote[2]{We have in our own literature a work on this subject
from the beginning of this century by C. \emph{Bastholm} (Historiske
Efterretninger til Kundskab om Mennesket i dets vilde og raa Tilstand. 4 vols.
Copenhagen 1803---4). Under the title of ``Descriptive Sociology''
\emph{Herbert Spencer} with several collaborators has begun to publish tabular
surveys of the cultural-historical conditions of the various peoples.} They
are like experiments that nature and history have made for us in a domain
where we \opage{64}cannot experiment; and yet in the preceding we have made
the relations still simpler than they are in reality; for as soon as man has
raised himself to the real beginning of a social life, he feels himself
related not merely to the authorities that rule in family, tribe and state,
but to a still higher authority, which lies beyond the bounds of these
societies, and in which at last those authorities themselves seek their
support: the religious authority. Only through this last support can the
principle of authority become absolute. As one of this principle's most
consistent defenders, Bonald, has said: ``Religion, which is the general bond
in every society, tightens preeminently the knot of political society; the
very word religion (religare) indicates sufficiently that it is the natural
and necessary bond of human societies, families and states. --- Religion
brings order into society, because it teaches men whence power and duties
stem. --- All power is created in God's image and stems from
God.''\footnote[1]{G. \emph{Brandes}: Reaktionen i Frankrig. S. 110; 118 f.}

If one regards as the essential thing in religion the feeling of dependence on
a higher order of things, one must in reality grant that those are right who
maintain that there are peoples without religion. Such a feeling always
presupposes a certain spiritual development, which is not found among human
beings who lack the capacity to form general ideas, or who have not yet
thought about \opage{65}whether the sun that rises in the morning is the same
one that sets in the evening.\footnote[1]{\emph{Lubbock}: Origin of
civilisation, p. 10; 475. (The languages of many savage peoples have names for
each particular color, each particular kind of tree, but none to denote color
or tree in general).} If on the other hand one demands of religion only belief
in spiritual beings,\footnote[2]{\emph{Tylor}: Primitive culture. I. p. 382.}
then one cannot deny that it is found in the crudest fetish-worshipper, who
does not even harbor fear of his fetish but merely regards it as a necessary
means and tool. A trace of a moral relation shows itself already within
fetishism, in that the Negro veils the fetish (as the Russian peasant does his
saint) when he does something he is ashamed of. Even where greater and
mightier beings are worshipped, fear is long the ruling feeling: one worships
only the evil beings; the good beings, one thinks, need no attention, since it
is their nature to be good. This dualism between good and evil beings is
therefore not of an ethical nature either, but springs from the utility or
harm that man has experienced. Thus the sun is in cold countries a good being,
in hot countries an evil being. It is anything but friendly feelings that men
at this stage harbor toward the gods they worship.

Only in idolatry proper, where anthropomorphism has made its way through, does
religion become real submission. Ethical ideas are transferred from political
society, where their \opage{66}first development has taken place, to the world
of the gods. The human race seems to have reached the standpoint of monarchy
in politics before it reached that of idolatry in religion. The deities now
gradually change from mere powers of nature into ethical powers, a transition
that can be followed especially within Greek mythology. And just as the belief
in gods changes, so too does the belief in immortality; men now no longer
believe merely in a continuation of existence indifferent in itself, but in a
retribution after death. The standard of this retribution is given with what
the people commend as virtues or vices.\footnote[1]{Cf. \emph{Tylor's} and
\emph{Lubbock's} above-mentioned works and \emph{Preller's} ``Griechische
Mythologie''.}

It marks a higher stage when the good, too, is recognized and worshipped. With
gratitude is combined admiration of the great and the exalted. The Fijians'
word for God (kalon) means in general whatever is great and wonderful; it is a
superlative expression. And at the highest stage of folk religion that we
know, the religious idea raises its object wholly above nature, in that
everything in the world is seen as brought forth by the will of the deity. The
feeling of dependence first becomes complete here, in that man stands, over
against omnipotence, as absolutely powerless. What the omnipotent divine will
commands is the highest law. A firmer and deeper foundation for the ethical
law seems not to be findable. Here authority has reached its highest
\opage{67}form, appears as absolute, unconditional. Authority in family and
state is only derived in relation to it. And although the divine authority has
for the believer clothed itself in the attributes of wisdom, love and justice,
so that he worships in it not merely the infinite power but also the ideal of
holiness, it is nevertheless the omnipotence of the will, and the fear and
trembling that the idea of it must awaken, that is the ultimate foundation.
The older theologians state this without reserve. It is the desire for eternal
blessedness that leads man to God; ``men,'' says Augustine, ``ought to worship
no god unless he can make them blessed; if blessedness were a goddess, she
would be the only one who ought to be worshipped.'' Obedience to God's will is
therefore, as Augustine likewise shows, the fundamental Christian
virtue.\footnote[1]{De civitate dei. lib. V. præf. --- XIV. c. 12.} The
consistent carrying through of this standpoint therefore cannot give up the
doctrine of the punishments of hell, for without them authority would be
weakened. They are the keystone of the edifice of absolute authority. ---

These features of the history of authority were meant only to illustrate, to
some extent, a principle that has intervened very mightily in the ethical
development of the human race. There was a time when, in revolutionary zeal,
one had eyes only for the less fortunate aspects of the principle of authority
and could not name it except in order to declaim against it. Now most, to be
sure, agree about its great historical mission, which \opage{68}is not yet
ended, and one is perhaps often inclined, out of counter-revolutionary zeal,
to regard its cultural-historical significance as a proof of its truth in
principle. The conduct of many modern apologists recalls, as \emph{Vacherot}
has aptly said, Scipio Africanus, who answered the charge of embezzlement with
the words: ``Let us go to the Capitol and thank the gods because I beat
Hannibal at Zama!'' The crowd followed him, but the accounts were hardly the
more correct for that. We cannot avoid settling the accounts. We were led to
authority as to the historical power that originally awakened the idea of a
higher law in man's breast, and on which this law has hitherto rested. Just as
necessary as it was for us to bring forward the concept of authority, just so
necessary is it now to examine whether it is necessarily bound up with the
concept of law, and which of these two concepts is, ethically regarded, the
essential one. Since it is a question of principle, we must take up authority
in its completed, its unconditional form, and this we have, as shown, in the
religious authority. The question of principle that we raise is the old point
of contention between theological and philosophical, Christian and humane
ethics.

The law presupposes an authority to which man bows. But with this it is not
yet given that the authority itself is one with the law. If I bow to the law
because it is the expression of the authority's will, then my motives have not
sprung from the idea of the law itself \opage{69}and its content, but must be
sought in the relation in which I, apart from the law, stand to the one who
commands, be this relation now one of fear, of piety or of reverence. Thus it
is only by a detour that I arrive at recognition of the law. If it were
another who uttered it, I would perhaps not obey it; but I would also obey a
contrary command if it went out from the same authority. The content of the
law thus becomes something accidental, and man is caught in the contradiction
that his most energetic action, that into which he puts everything, stands in
a purely arbitrary relation to his own nature. A third power always thrusts
itself in between him and that for which he works, and an unconditional power
at that, before which everything else becomes nothing. The unconditional
authority is a ``jealous God,'' who tolerates nothing beside himself. And even
if the authority is not absolute, --- when it cannot make itself one with the
content of the law, it steps disturbingly in between man and his task, hinders
the direct, inward relation.

In the (original, consistent) Christian ethics this comes clearly to light. It
is said in general that the commandment of love is the chief law of Christian
ethics. In this commandment, great and exalted as it is, there lies in itself
nothing that points beyond the humane. To love others, apart from all
differences, and to love them unconditionally, is a commandment that Buddha
had already uttered. The Stoic philosophers had likewise earlier set up love
\opage{70}for the whole human race (caritas generis humani) as the fundamental
ethical law, although they were not able to procure for it so deep and inward
a recognition, even among the simple, as the Gospel, which precisely thereby
proved its historical mission. Here, where we have to do only with principles,
we must disregard the practical significance they may have under certain
historical conditions. It then turns out on closer consideration that it is
not for men's own sake that we are to love them, but for God's sake. Love is
not the fundamental Christian virtue, but obedience, submission to the divine
will. Without this, Christianity would not be a positive religion. Positive
religion means religion grounded on authority, and obedience is the
fundamental relation toward authority. In Christianity this finds its more
precise expression in love's being subordinated to faith. The apostle Paul
indeed teaches that love is greater than faith, but Christianity recognizes
only the love that is grounded on faith, faith in the supernatural authority.
Therefore the virtues of the heathen are, in the eyes of earnest Christians,
only splendid vices. In its commandment of love Christianity, like Buddha and
the Stoics before it, pointed beyond the outward differences that divided the
ancient world: there shall no longer be any difference between Greek and
barbarian, between bond and free; but it has itself, in return, established an
opposition that is surely as strong as any of those: that between believers
and unbelievers, the blessed and the
% continues on p. 71
% continues from p. 70
\opage{71}Damned. Faith and love therefore work against each
other,\footnote[1]{See \emph{Ludwig Feuerbach}: Wesen des Christenthums (Cf.
my book ``Filosofien i Tyskland efter Hegel'' p. 65---67).} and the
commandment of love first comes into its full right when one goes beyond the
bounds of faith. But this Christianity, as a \emph{positive} religion, cannot
do. Therefore the commandment of love has first attained its completion in the
modern principle of toleration, which treats the differences of faith in the
same way as Christianity treated the outward and national differences. Spinoza
is the first who has drawn the consequences of the commandment of love and
carried through in practice the proposition that love is greater than faith.
The humane element in Christianity has hereby been freed from the bounds that
hampered it.

It is not only love that acquires a merely indirect significance under the
presupposition of the principle of authority; the same holds of every humane
and cultural-historical relation. Marriage, on the strict Christian view, can
only become the lesser of two evils. Paul advises the slave against seeking to
obtain freedom, and Luther says in his blunt language: ``Serfdom is not
contrary to the Christian essence, and whoever says so, he lies.'' Before the
one infinite object that is to fill the minds of all believers, everything
that has only human significance must be pushed aside.

One will perhaps object that we here push the \opage{72}matter to an extreme
and conceive authority as the purely arbitrary. ``Certainly,'' one will say,
``we accept the good as good only because it is God's will; but since the good
is one with God's essence, he wills only what is good, and so what becomes the
content of the law does not rest on arbitrariness and chance.'' But this is
only to push the whole question further back, unless one ascribes to oneself
so intimate an acquaintance with the essence of the Godhead that one knows the
reasons that determine its will. And if one knows these reasons, one can see
their correctness; why then take the detour of appealing to the supernatural
will, instead of applying them directly? As soon as authority does not wholly
subdue us by an unconditional command of power, it cannot avoid being called
to account; but an account presupposes that there is something that stands
above authority, and by which the latter can be judged.

Although \emph{Martensen} in his ``kristelige Ethik'' % misprint in the print: the closing quotation mark is an opening «
energetically maintains the necessity of the supernatural, unconditional
authority as the foundation of the ethical, he has nevertheless clearly seen
that one must establish the justification of authority. He solves the question
by saying: ``The right of authority is grounded in the ethical, to which power
is subordinate'' (p. 449). This proposition we could subscribe to entirely.
But what, then, does Martensen here understand by ``the ethical''? If it is to
be of the same nature as the ethical whose explanation we sought, then
according to Martensen's teaching we must of course again seek a higher
authority, whose will \opage{73}can give it validity --- a power, then, that
can obligate God! And if the proposition that the ethical presupposes an
obligating power standing above the one who is obligated is to hold only for
the humanly ethical, not for the ethical in the essence of the Godhead, then
this latter ``ethical'' becomes for us an empty word, a sign that one has
applied the concept where it no longer has any validity, let alone that it can
serve for any explanation whatsoever. Beyond this dilemma theological ethics
cannot get.

Where authority is absolute, i.e., where it is raised above all grounding, it
must make itself known by a miracle. For the theological conception,
therefore, the connection of nature is broken at the point where the ethical
begins. ``The ethical `Ought','' says Martensen, ``remains in all eternity
inexplicable from nature. It comes not from below, but from above'' (p. 441).
And later he maintains that the majestic: ``Thou shalt!'' does not even find
its explanation in our ideal nature (p. 456). But for Martensen spiritual life
as a whole is ``a miracle for nature, a transcendent for the whole concept of
nature'' (p. 16). What we must object here is that theology constantly
confuses the (as yet, or in itself) inexplicable with the miracle. The miracle
is a fiat, an asylum of ignorance (asylum ignorantiæ), in which theology takes
refuge in order to fill the holes it believes it has discovered in scientific
knowledge. Humane ethics, which must apply the same scientific procedure that
holds everywhere else in the world of thought and of nature, does not
\opage{74}imagine that it can solve all riddles; it seeks to get as far as it
can by the natural ways; the rest it leaves undecided. But it cannot proceed
from so narrow a concept of nature as theology does. Although consciousness
cannot be derived from the unconscious life of nature, it can nevertheless be
shown that, as regards the conditions of its origin and development, it stands
in the closest connection with it and cannot be understood without this
connection. The psychological and ethical laws are therefore laws of nature,
just as spiritual life as a whole is a link in the infinite life of nature
itself and emerges at a certain definite stage of development of the latter.
To the popular conception it does of course appear as if certain spiritual
states, thoughts and feelings had so exalted a character that they could not
possibly be brought into connection with the ordinary psychological laws. Thus
the motions of the heavenly bodies stood for the ancients as divine and ideal,
so that it was a presumption to explain them in a natural way; and yet in
modern times they have found their mechanical explanation. However far
psychology and ethics lag behind astronomy, they nevertheless rest on the same
principle; they too strive after a ``mécanique céleste'', a natural
explanation of ``the heavenly things'' in the domain of spiritual life. Faith,
worship and love will find their psychological and historical explanation,
just as the more elementary phenomena of consciousness have already found
theirs.

It is the task of historical science to \opage{75}describe for us the origin
and development of the authorities. These are to be regarded as natural
beings, which unfold according to definite laws. Historical development and
its laws stand above all authorities and comprise them all. When the
conditions of an authority's validity are past, it gives way, as history
shows, to another. History is the battleground of the authorities. But
precisely in their fighting lies the fact that they are, after all, finite
beings. They are like the gods of Olympus, who descended onto the battlefield
and could be wounded by men.

Herein we find the possibility of recognizing the great significance of the
authorities for the ethical and yet maintaining that the principle of
authority in itself is not the full ground of the ethical. They are the
educating powers in the history of the race. But he who educates another
strives, of course, precisely to free him, to bring him to stand on his own
feet and see with his own eyes. He strives to make himself superfluous, to
disappear himself, so that the disciple can come face to face with the matter
itself. This resignation is demanded of every educator, but the authorities in
every domain forget it all too often. Precisely hereby they prove their
finitude. For that whereby the authority differs from the law that is to be
upheld, the end that is to be attained, is the egoism, the self-will, that it
asserts at the expense of the cause that was to be furthered. With this it
agrees well that authority always in the last instance comes to appeal to
egoistic motives, above all to fear. \opage{76}Ethical validity it has only as
the bearer of a content that can prove its justification quite apart from the
authority, and only thereby does it become an object of admiration, reverence
and love; if, on the other hand, it acts beyond the justification this content
gives, it sinks back to its most primitive forms. As long as the authority
stands as a third between the agent on the one side and the law and end of the
action on the other, so long is action only indirectly ethical, because other
motives rule than the very recognition of the law and reverence for it. Ethics
thus reduces all authorities to being relative instead of absolute, means
instead of ends. Authority exists for man's sake, not man for authority's.
When the personality has developed so that the recognized law is his only, or
at any rate his deepest, motive, then there is in him an authority before
which all outward powers are nothing. The center of gravity is shifted from
the outer to the inner, --- the great goal of all education!

It is the task of modern times to transform the authorities from absolute to
relative. Instead of the weight under which life groaned so long as a
supernatural authority, partly immediately, partly through its representatives
and servants, embraced everything, great and small, with its commandments and
its prohibitions, modern times work to bring about a freer and higher
development by making the heavy thing the foundation instead of the roof and
spire. For this the Reformation and the Revolution fought, each in its
\opage{77}way. Authority is a condition of life for social life, and thereby
also for the ethical life, but life first reaches its higher development when
authority subordinates itself. The conviction will make its way that the
ethical ideas have their validity and their worth in themselves, through their
intimate connection with the essence and conditions of human life, although
historically they develop under the shelter of authority. By working to make
the latter relative instead of absolute, one therefore works in the service of
the ethical. For the consciousness of most people, to be sure, the ethical
ideas have entered into so close a connection with the idea of the absolute
authority that an attack on the latter also seems to them to be an attack on
the ethical; hence the constant accusation against free inquiry that it leads
to unethical consequences. Just as the pagans in Alexandria expected that
heaven and earth would fall back into the old chaos when the Christians
overthrew the statue of Serapis,\footnote[1]{\emph{Gibbon}: The decline and
fall of the Roman Empire. Chap. 28.} so many now think that human life will,
in ethical respects, dissolve into a chaos when the old dogmas are toppled
from the place they have hitherto occupied as ruling powers. For him who has a
freer view and a bolder faith the matter appears otherwise. Without fantastic
illusions, and without being frightened by the passing excesses that never
fail to appear where great and new ideas are breaking through, he will leave
it to \opage{78}the quiet but sure course of historical development to bring
harmony between the practical forms of life and the new ideas. We have seen in
an earlier connection of how essential importance it is for the individual
human being to be in harmony with his own highest knowledge. The same must
hold for society as a whole, for the whole race: if an idea is thus really
grounded in human nature, it will also work its way forward to practical
recognition, however little we may be able to survey all the intermediate
stages, the ways and detours, that must first be traversed here. Here, then,
\emph{Pascal} is right when he says: ``Let us work at thinking rightly; that
is the principle of morality!'' That is at least the contribution that
theoretical ethics ought to make to the development of the ethical life.

Only those authorities, then, are ethically justified which see their own
relative significance. The supernatural authority must always answer every
call to subordinate itself in the same way as its most consistent
representative, the Roman Church, answers the call to recognize the progress
of history: ``Impossible!'' (Non possumus). The essence of the humane
authorities, on the other hand, is not violated by such a call. Authority is
here essentially an expression of tradition, which the one generation
transmits to the next, giving it along on the way as an indispensable support,
until it can stand on its own feet. Only thereby does continuity arise in the
human race, and only thereby do progress and development become possible, that
each new generation does not begin \opage{79}afresh, but stands on the
shoulders of the preceding ones. The clash between the older and the younger
generation rests partly on the inherited forms not always yielding after their
proper role is played out, partly also on the unbridled and impatient eruption
of the new forms. But the very hard struggles that can arise from this have
their great significance: freedom must be won through struggle; the
recognition of the inner, higher law cannot be grafted on from outside. The
resistance that even the relative authorities are thus inclined to offer to
their own subordination serves precisely for a more inward, freer and more
vigorous development of the new forms. Yet it must always stand as an
essential task to make this transition as short and as easy as possible, and
this will become possible the more the conviction spreads that ``the right of
authority is grounded in the ethical'': the more the authorities really act as
educating powers, and the more inwardly and deeply their significance is
therefore also recognized by all. The ethical law is not only to issue from
authority, but also to be applied to it. An absolute authority is raised above
all laws, even above those it has itself given; but the relative authorities
one can measure by their own yardstick.

As a member of the race, each finds himself set to solve his particular task
as it emerges from all the historical conditions. He is a child of the past
and helps to create the future. He cannot see his dependence on what has gone
before without at the same time seeing that what is to come is in part
determined by his conduct. However limited \opage{80}his horizon and his
sphere of activity may be, he is never excluded from regarding himself as a
fellow worker in humanity's great society, advancing through the ages. The
consciousness of this, under whatever form it may show itself, is one with the
feeling of obligation, with the recognition of a higher law. Here we have the
ethical in its simplest and purest form. \emph{Comte} has beautifully and
aptly shown how that consciousness arises within family life and constantly
points back to it. Through the relation of piety toward the parents the child
enters into relation with the past of the race and at once receives a deep
impression of how closely its whole existence is bound up with it. In the
relation of brothers, and later in the relation to fellow workers and
contemporaries in general, the individual is led into the whole great circle
of solidary interests, where he too is to contribute his mite. The connection
does not stand out here so clearly as within the family, but the further
culture advances, the more closely bound up all human endeavors nevertheless
prove to be, so that progress in one direction is in the long run impossible
without corresponding progress in other directions. The relation of father,
finally, forms the conclusion and points toward the future. What appears in
its simplest form within the family can be extended and developed to comprise
the whole race. There is one great circle of tasks that is essentially
conditioned by the relation to the past; it comprises everything that serves
to preserve and secure human life, and the organization, the fixed forms,
without which life can here no more \opage{81}than anywhere else subsist.
Another great circle of tasks points toward the future; to it belongs
everything whereby life can receive higher and wider development and
enrichment. Order and progress presuppose each other: order is the condition
of progress, and progress is the end of order.\footnote[1]{Politique positive
I. 95; 105.} In both circles of tasks personal life finds its content and its
nourishment. The harmonious society of free personalities, which in this
connection too proves to be the highest ethical idea, is not a mere ideal, but
is relatively and partly realized in every ethically developed personality. If
the goal always lay in the future, ethical action would take on the character
of a Tantalean chase after something that always eludes us. But the kingdom
for whose coming we are to work has already come where the work is done
rightly, whether the work in the particular case is concentrated chiefly on
the development of the individual's own personal life in its various
directions, or has a more outward-going character. The connection with the
great and common end of the human race hallows and ennobles even the humblest
deed, when it is performed with faithfulness.

It is not merely abstract laws and reflections that meet thought in the
contemplation of the development of human life. The great figures of history,
in whom the human tasks and ends appear as though personified, stand as
\opage{82}models, as types of what human life can be. Some of these figures
have, even in the form in which they have been handed down to us, such a stamp
of perfection that they seem raised above the struggles and difficulties
through which it is the lot of most to work their way forward. Myth and legend
have in many cases been inspired by the same thought that \emph{Schiller}
expresses in one of his epigrams, that noble natures count by what they
\emph{are}, lower natures by what they \emph{do}. It is a thought that has its
great justification, since, as was just remarked, the way and the goal,
ethically considered, cannot fall entirely outside each other. But ``the
church militant'' needs models who have \emph{really} fought themselves,
for whom the struggle has been more than a divine game. The model can
therefore no more be absolute than authority can: that would be the same as
making it useless for relative beings. What we are to learn from it is
precisely what human life can be made into by the exertion of man's own
powers.

Life of great men all remind us,\\
We can make our life sublime. (Longfellow.)

However mighty an influence the thought of progressing humanity may have,
especially when it is seen represented in the outstanding personalities, the
consideration of the ethical nevertheless leads us further. Development did
not first begin with the appearance of man. The best-grounded
\opage{83}hypotheses\footnote[1]{Cf. my lecture ``Filosofien og Darwinismen''
(``Nær og Fjern''. 1874. Nr. 93---94).} lead us to see the origin of the human
race itself as the result of a long organic and physical process of
development. Through a series of stages, under the most diverse forms, the
life of nature rises from the dim nebular masses to the highest spiritual
clarity. Everywhere we can find the same fundamental laws, the progressive
formation of smaller wholes within the larger whole.\footnote[2]{Cf.
\emph{Herbert Spencer}: First Principles.} Every individual form, in the
organic as in the inorganic, in the human as in the animal world, marks a
peculiar stage in the great development of life. We do not explain the higher
from the lower because we see the former emerge from the latter; on the
contrary, we first rightly understand the life of the world when we see what
it is able to achieve at its summits, in its most individual forms. Man,
precisely according to the hypothesis of evolution, is no stranger in the
world, but has his roots deep in its innermost essence; and what stands for
man as necessary, as true and as good, is just as much a reality as the
outward, material masses. By working for the realization of his ideals, by
developing and perfecting human life, he plays the role that is to be played
at the place where he is set. He stands as the final end of a long
development, and his natural task is to fill this place in such a way that he
not only maintains \opage{84}it, but also uses it to carry the development
further. His struggle for existence becomes a struggle for the dignity of man.
It is the spiritual life, the life of thought, of feeling and of will, whereby
man distinguishes himself from the lower stages of life, and which he must
guard and develop if he does not want to sink back from the place he has won.
All the individual starting points that we have brought forward earlier now
acquire their proper significance by being fitted into the great connection in
which human life stands. The \emph{impulse of self-preservation}, which by
concentrating one-sidedly on the single individual becomes egoism, acquires
its justification when it is seen that the individual's struggle to subsist is
a condition of the subsistence and development of the whole race. The tasks
that set the individual interests in motion are at the same time those in the
handling of which the race exercises and develops its powers. \emph{Sympathy}
becomes the uniting bond between the race's simultaneously and successively
living individuals; it rises above mere instinct when it is purified by the
thought of the great common task on which all work, at all times and under all
conditions. \emph{Reason} now becomes more than a purely formal capacity; it
acquires its living and full content through the contemplation of humanity's
course of life, the laws and conditions under which it takes place, the tasks
and duties that follow from these under given conditions.

Although human life thus has its ground in the life of the world, it is
nevertheless a vanishing quantity in relation to the latter. The times are
past when one could \opage{85}rest content with the assumption that man was
the center of the life of the world, and everything existed for his sake. The
more thought has penetrated nature and learned to know its laws, and the more
man, by the help of this insight, is able to use the forces of nature in his
service, the more has the knowledge awakened at the same time that he himself,
that the whole of human life, is only a vanishing element in relation to the
great whole. The infinity of the life of the world and the unshakable laws
according to which it unfolds convince him of his finitude and limitation. The
feeling of dependence, in which the essence of religion has rightly been
sought, will thus naturally arise also in him who does not live in faith in
supernatural powers. But all fear and egoistic interest are excluded by the
clear knowledge that what is to become reality rests not on the wishes and
feelings of men, but on a definite order of things, expressed in fixed laws.
He who strives for knowledge of himself and the conditions of his existence
will therefore, as a ``joyful link in the chain of the world
order'',\footnote[1]{\emph{Dante}: Paradiset, 3rd Canto, 54th Verse.}
% Danish: «frydfuldt Led i Verdensordenens Kjæde»
in bold resignation leave to the course of nature what, and how much, of what
he desires and hopes for is to be realized for him. When we hold fast that the
highest we know, the ideal human ends and laws, has developed through this
order of things, the feeling of dependence becomes reverence; the world order
becomes not merely physical, but ethical. However \opage{86}vanishing human
life may be in relation to the life of the world, the highest goods we know
have their place within it, and their worth is not diminished by our being
unable to make sun, moon and stars stand still or change their course in order
to serve them. We owe life and existence far too much to make such demands. In
our deep dependence on nature we nevertheless feel gratitude and reverence for
the great connection that has also let the fruits of spiritual life ripen.

The ethical thus stands in close connection with the religious. In both
domains the feeling of reverence plays a leading role. The difference rests on
this, that within the ethical, where man stands as acting, this feeling is
attached to his own highest knowledge of ends and tasks, but in the religious
he goes further back and sees his action, as well as his whole being, all his
wishes and ends, as a link in the whole of human and natural life and subject
to its conditions. So long as the feeling has not been extinguished in
humanity that the individual as well as the race are links in the infinite
whole, issue from it and are conditioned by it, so long will the ethical and
the religious also continue to assert themselves under various forms. We stand
here before something that has its firm ground in man's nature and conditions,
and that does not stand and fall with the commands of arbitrary powers.

% ===== THE ETHICAL LAW AND PROGRESS (pp. 87--107) =====
\chap{--- The Ethical Law and Progress.}
\opage{87}The adherents of the belief in authority praise it as one of the
most important advantages of their principle that one learns directly from
the source of truth what the good and the right are, so that the decision
on this is not left to the changing opinions of times and individuals.
This would certainly be a great advantage if the beings who were to act
were incapable of striving themselves after a conviction in these
questions, and if the chief weight were to be laid on a certain sum and
kind of actions and not on the free disposition from which the actions
spring. But when one sees man's highest powers stirring precisely in his
striving after knowledge of his nature and his position in nature, one
will grant Sir \emph{William Hamilton} that a living error is better than a
dead truth, and the living striving after ethical conviction itself
becomes an important element in the ethical, without which the kingdom of
harmonious \opage{88}personalities cannot be built up. The man who stands
and listens for commands from outside before he can decide has not yet
attained true personality. It will moreover always prove to be an illusion
when one thinks that life is in reality governed by the commandment of the
supernatural authority. This stands too far off to be able to make its way
through in face of the actual conditions, where so many other forces ---
noble and ignoble --- are stirring. ``Not one Christian in a thousand,''
says \emph{Stuart Mill} (Om Friheden, ch. 2), ``lets the principles and
precepts contained in the New Testament be a guide and test for his
individual conduct. The standard by which he directs himself is what is
custom and usage in his nation, his class of society, his religious
profession. He thus has, on the one hand, a collection of ethical
principles which he believes an infallible wisdom has given him as rules
by which he is to direct his conduct; but on the other hand he has a
collection of everyday opinions and directions which on the whole are a
compromise between the Christian faith and the interests and promptings of
actual life, in that they to a certain degree agree with some of those
principles, do not agree so much with others of them, and finally stand in
direct conflict with certain of them. To the first of these guides he
gives his homage; by the second he in reality directs himself.'' But it is
not only tradition and habit that can thus hinder the \opage{89}influence
of the unconditional authority; the free, independent conviction about the
demands of humanity also plays more and more a considerable role, to which
the whole social development of modern times sufficiently testifies. Here,
then, is visible proof that a firm insight into the ethical can really be
acquired in a free and human way. To overlook this source in order to seek
another would be to cross the stream to fetch water.

Only then can one rightly praise the principle of authority for its
certainty and clarity, when it is possible in every single case to get an
answer to one's questions. The Catholic Church affords its adherents this
advantage. The Catholic can in all cases of doubt turn to his confessor,
who stands as the representative of the Church and thus, in the last
instance, of God, as a pipe through which the flow from the source of
absolute truth is led to him, the individual. Protestantism forgoes this
certainty and clarity when it refers each individual to his inner
conviction, and it takes up a wavering and untenable position when it
would remedy this by ascribing to the Church a relative infallibility
(\emph{Martensen's} ``Katholicisme og Protestantisme''). A relative
infallibility is not infallibility. For the Protestant, conscience in
reality becomes the only guide, and he can follow the Church only where it
agrees with conscience. Protestantism takes a step forward with respect to
ethical emancipation, but remains standing halfway. This comes out
characteristically in the oldest Protestant moral theologians,
\opage{90}who taught that conscience ``in respect of power and authority
stands midway between God and man: under God, but above man.'' Here the
possibility is still hinted at of an authority that could stand above
conscience, and it is this possibility that separates Protestant ethics
from humane ethics.

If it stands fast that man can be bound only in conscience, by free
recognition and feeling of the right, then conscience itself is the
highest authority, and subordinate to no other whatsoever. Conscience, after
all, means the highest degree of clarity and inwardness to which man's
ethical thinking and feeling has attained by every possible exertion. Since
it is now man himself who is to act, and man as he is under certain given
conditions, he can seek the standard for his action only in the highest
understanding of the nature and significance of the action that it is
possible for him to gain. We must therefore grant \emph{Johann Gottlieb
Fichte} his bold proposition: ``Conscience never errs and cannot err.'' It
holds for man in the moment of action. If the better insight comes before
the action, it will if possible prevent it; and if it comes after the
action, it comes as a judge who must nevertheless acknowledge that the man
did right in acting according to his best conviction. It is a duty to have
a conscience; man must strive after a clear and definite conviction about
the right. At every given stage he must decide according to the insight
that \opage{91}stands at his disposal, while he carefully seeks to remedy
his limited knowledge by anticipating and imagining how the action would
look from a higher point of view. But since this higher point of view
itself is always modified by the actual standpoint and circumstances of the
one acting, which he is able to look away from only in part, the
possibility is never excluded that the action may later, under other
conditions, look different to him. Nevertheless, if he really acted from
conviction, he will see that the later, clearer and more correct insight
itself became possible only through his having applied his imperfect
insight at the earlier stage as well as he was able. Both stages thus lie
on one and the same line, mark stages of the development toward one and
the same goal.

What holds for the individual holds for the human race as a whole. The
various human stages of development mark just as many various stages in
the knowledge of the ethical law. One cannot maintain that unity prevails
among men in ethical principles; experience and history refute that. Unity
is not the starting point but the final goal. But what appears everywhere,
as soon as a certain human development has been reached, is the relation
of obligation itself. This can be present with greater or less clarity and
inwardness, depth and range. What in one place is drawn under this
relation is regarded in another place as indifferent. As soon as a higher
stage has been reached, judgment is passed on the \opage{92}lower, through
which it has nevertheless itself become possible. In so far, history is
the judgment in the life of the race as well as in that of the individual.
An absolute truth we grasp nowhere. It is our lot to work at the solution
of an infinite riddle, which presents itself under ever new forms as long
as life lasts.

We are hereby led to the question whether one has the right to assume an
advance in the ethical respect in the human race. This assumption, to
which so much in our preceding investigations must lead us, has been
rejected, by holding chiefly to the inner, subjective side of the
ethical. Thus the famous naturalist \emph{Wallace}, who knows the Malay race
thoroughly, declares that the savages fulfil their simple moral law at
least as well as the civilized Europeans, indeed that the mass of our
population has not got beyond the moral law of the savages, but in certain
respects has even sunk below it. One has gone a step further, in rejecting
the thought of a properly ethical progress as containing an injustice,
since the later generations, which began their course at a higher ethical
stage of development, would then be more fortunately placed than the
earlier ones. Since in the ethical what essentially matters is what one
makes of the given conditions, it makes no difference at which stage one
begins. Under the simplest and the most developed conditions of culture,
morality can thus be the same. On this way of looking at things, which in
our time has been maintained by \emph{Bouillier}, the ethical is the same
\opage{93}at all times: numerator and denominator may vary, but the
fundamental ratio remains the same. (1/2 = 2/4 = 3/6 = 4/8 etc.).

One here assumes that the ethical relation is present even at the lowest
human stage of development. But the ethical relation presupposes
definite conditions, which cannot always have been present. It cannot be
definitely shown before society and the relation of authority have
formed, and before man's spiritual capacities have reached a certain
strength and a certain range. Originally man distinguishes between good
and evil only as between that which brings him joy and that which brings
pain. Thus a savage explained that when someone took his wife from him, it
was evil, but when he took another man's wife, it was good. And after the
ethical relation has come into force as recognition of a higher law and
action from reverence for it, it will take on a very different character
under the different historical circumstances. The ethical is not something
isolated in man's nature, does not stand indifferent to his development or
decline in all other respects, but enters into the most intimate and lively
interaction with everything else that stirs in man. The range, the purity
and freedom with which the ethical asserts itself will always to a certain
degree depend on the stage of culture and spiritual development that man
has attained. With a clearer understanding of the goal and the conditions,
the obligation itself must also take on a higher character. Even as
\opage{94}regards the fundamental ethical relation, we must therefore
assume an advance. But an entirely different question from this is whether
the later generations thereby gain any advantage over the earlier ones.
The conditions under which development begins for the several generations
and individuals are not in their own power. We naturally occupy a higher
standpoint than those on whose shoulders we stand, but that we do so we owe
precisely to them. And the more the connection between the generations
comes forward into consciousness, the more, in particular, the educating
authority of the preceding generation over the one growing up is seen in
its whole significance, the more will piety also be awakened toward all
those who, each in his own way, have fought and worked to set us where we
now stand.

The necessity of assuming an ethical advance is seen still more clearly
when we look to the objective side of the matter, to the forms of life
and institutions in which the ethical procures itself expression. The
freer and more human way itself in which the conditions of life are
gradually ordered and arranged for men must surely be of higher ethical
worth than the crude and primitive arrangements. It is accumulated labor,
which, like all capital, can certainly be both used and misused, but which
nevertheless represents a fund at the disposal of the human race.

From the ethical point of view we regard the history of the human race as
the history of the development of the harmonious society, the progressive
realization of the kingdom of humanity. Within this development one can
distinguish three stages; in the first \opage{95}might rules, in the second
right, in the third duty. These stages can be shown for the race as a
whole as well as in the ethical development of each single individual and
within the several humane and social spheres.

No education takes place without compulsion and without the use of
punishment. Even the most philanthropic educator must in practice use
compulsion, in order to prevent his ward from actions that would be
absolutely and immediately harmful to himself or others, and the ward will
often for the time being receive no other impression of this than of any
use whatever of superior force. In primitive social life, where the
philanthropic expresses itself only weakly, might rules almost completely,
guided by the impulse of self-preservation of the mighty. All labor is
loaded onto the weak, onto women and children. The vanquished enemy is
treated like any other booty. Only within the tribe does one recognize the
right of other individuals to exist and to satisfy their interests, and
this recognition itself has been extorted by might. Right rests originally
on recognition of the given relation of might. Within the family, where
the relation of right does not assert itself, might continues to rule, as
is seen in particular in the older Roman law, which subjects wife and
children completely to the authority of the father of the family. When
compassion for the weak, which is developed and nourished precisely within
the family, leads to the defenseless enemy no longer being killed, he is
made a slave and as \opage{96}such still has no more share than wife and
children in the relation of right. The master still stands toward him in
the state of nature; war smoulders behind peace, as is seen from such laws
as the one in Rome which ordained that all of a man's slaves were to be
killed when he had been murdered, if they had at the time been so near
that they could hear his cries.\footnote[1]{Cf. \emph{Tacitus}: Annales
XIV, 42---44 (cf. XIII, 32.)} This law sprang, as \emph{Montesquieu} aptly
remarks,\footnote[2]{L'esprit des lois XI, 16.} from the law of war, only
that here the enemies were in the interior of society. How the slaves were
treated depended partly on the master's disposition, partly on external
circumstances. Of Cato the Elder it is told that he let his slaves judge
one another among themselves, thus establishing a kind of relation of
right among them of his own good will. But when suspicion constantly stood
in the background, when the master felt himself the one among the many,
full security could never come about. The humanity with which the slaves
were nevertheless on the whole treated among the Greeks and Romans
certainly had its ground not only in these peoples' humane and liberal
way of thinking, but also in the difficulty of procuring new slaves, just
as, conversely, the torments of the Negro slaves sprang in great part from
the ease with which they could be replaced, and which made the masters find
their account in letting them work themselves to death and then buying new
ones.\footnote[3]{\emph{Stuart Mill}: Principles of political economy. II,
5,1.}

\opage{97}The importance of the relation of might being replaced by the
relation of right rests on right's opening the way to a grounding. The
first element in the nature of right is certainly the sanction on the part
of the power of society, which undertakes to uphold everyone in the
possession and enjoyment of what he is and has. But with this a regulation,
an evening-out of accidents and inequalities, is necessarily combined. A
general standard and definite principles are sought for the decision of
the single cases; but this must in turn lead to reasons being given for
the single provisions. The old Roman jurists therefore sought to prove the
lawfulness of slavery.\footnote[1]{\emph{Hugo Grotius}: De jure belli et
pacis III, 7, 5; 11, 16.} They started from the premise that it is
permitted to kill one's adversary in war; if I nevertheless spare him after
having destroyed his power of resistance, he becomes my property. They
derived the word servus (slave) from servare (preserve), and in this way
presented slavery as a progress in humanity, which it in reality also was
in relation to cannibalism and the bestial raging against the disarmed
enemy. This grounding was still recognized in modern times by the first
philosophers of law (\emph{Grotius}, \emph{Hobbes}, \emph{Locke}). It was
reserved for the 18th century to deal this institution the decisive blow.
\emph{Montesquieu} overthrew the old grounding by the simple objection
that over a disarmed enemy, who can no longer do any harm, one reasonably
has no \opage{98}other right than that of securing oneself against renewed
danger.\footnote[1]{``The 17th century had no doubt with regard to
slavery. Bossuet accepted it without hesitation as a fact authorized by
Scripture. Locke did indeed combat it, but without originality and force,
and he even retained it in certain cases (criminals and prisoners of war).
The only discussion of it before Montesquieu that deserves to be
remembered is Bodin's in the 16th century. He was the only one in that
century who had raised his voice against slavery. As for the Middle Ages,
there was unanimity with regard to this institution, which Christianity
seemed to have destroyed. One must go back to the Church Fathers and the
Stoics to find so lively a protest as that of the 18th century. But the
Church Fathers, who relied only on the religious equality of men, and who
in the name of human right permitted what they rejected in the name of
mystical and Christian right, had not cut the root of this pernicious
evil. Hence it came about that, despite all the mitigations of slavery,
which had been transformed into serfdom, everything was ready for a
renewal of this evil, which had indeed been diminished but not destroyed,
when the discovery of America and of men with black skin gave occasion to
greed, superstition and ignorance. The voice of the doctors and the
theologians was not raised, with the exception of Las Cacas's, against
this attack on the right of man. One must therefore say that it is only
the 18th century that has dealt slavery a mortal blow; it is Montesquieu
who has had this courage and this honor; enlightened by him, by J. J.
Rousseau and others after them, the nations of Europe have resolved to
free themselves from this stain.'' \emph{Paul Janet}: Histoire de la
science politique. 2e. éd. II. p. 504---505.} % misprint in the print: "Las Cacas's" for Casas

By carrying through the line of thought that had thus been set in motion,
and through the enthusiasm for \opage{99}humanity and the respect for
individuality that were thereby awakened, one arrived in the last century
at the idea of universal and innate rights of man, which were formulated
variously by the various parties and in the various constitutions, and
which \emph{Kant} in his doctrine of right reduced to a single one: ``Freedom
(independence of every other's compelling arbitrariness), in so far as it
can coexist with every other's freedom according to a universal law, is
the only, original right which belongs to every man by virtue of his human
nature.'' It is, as one sees, the same principle that \emph{Stuart Mill}
in our day has so energetically and urgently enjoined in his book ``om
Friheden.'' Kant saw the significance of this principle expressed in the
fact that every restriction of personal freedom must be defended, whereas
freedom in itself needs no grounding. But precisely for that reason this
principle leads us beyond the proper sphere of right. For in so far as any
meaning at all can be connected with words such as ``innate right'' or
``right of man,'' it is an ethical rather than a juridical meaning. That
everyone is to be treated as a free personality is an ethical demand which
cannot in this form be fixed as an external law. But in addition, a purely
formal conception of this principle would precisely hinder its being
carried through. For it starts from the premise that everyone really is a
free personality, fit to have his center in himself --- but this is very
far from being the case. Freedom cannot be ``innate,'' it must be acquired.
\opage{100}By hastily recognizing men as free, one will precisely hinder
them from becoming so. It was also soon seen, when the intoxication of
enthusiasm over the ``rights of man'' was past, that it is of little use to
give man rank as a free being when one does not help him really to become
one. Here lies the social problem of our day, which in its manifold forms
touches all spheres, from religion, science and school to property,
industry and machinery. It is so comprehensive precisely because it does
not concern any single side of human life, but the full development of
personal life in all. Never before has the consciousness of what is at
stake been so lively.

We are hereby led to the conviction that the concept of right must be
subordinated to the concept of duty. Right consists essentially in
society's sanction of the individual's position of power. Right can
separate individuals and can cause them to regard society as a means to
their private ends. During the struggle against despotism a firm point of
support was found in the individual's absolute right of property. Property
is an atom! --- exclaimed the elder \emph{Pitt} in Parliament. But we know of no
absolute atoms. No atom is released from the general laws of the whole
nexus within which it finds itself. And just as the right to property
presupposes the sanction of the power of society, so property itself must,
from an ethical point of view, be seen as grounding a relation of duty.
\emph{Auguste Comte} has earnestly combated the modern distinction between
\opage{101}the private and the public; in his view everyone ought to regard
himself as a public functionary. Through property the individual, as
through his external position in general, has been assigned a sphere and
means for work in the service of the tasks of the history of civilization.
The ends of the history of civilization comprise everything that
contributes to the development and advancement of mankind in inner and
outer respects. Life and activity in the family, in civil society and in
the state, in the world of art and of science, of industry and of
commerce, serve this end. And by working in these spheres the individual
lives his way into the feeling that he is a member of a great organism, a
citizen in a kingdom of humanity that works its way forward in the course
of the ages. His proper self he then finds in the task that falls to him
within this development. The validity and significance of the personality
are conditioned by the value of the content and the interests in which it
lives. The one individual can assert itself and demand recognition from
the other only by establishing the significance of what it lives in and
for. The true personality is developed by work for ends and tasks set by
the actual world. Reality is more than a stage on which the individuals
are to present themselves in order to show their virtuosity. When the
monks in the Egyptian desert spent their time watering dry sticks that
were planted in the sand, this could be a lesson in ascetic self-surrender
and serve as a symbol of the insignificance of human activity in face of
the \opage{102}supernatural world. But the vital nerve of humane ethics is,
as we have seen from several sides, the intimate connection between the
individual and the race, and therefore also between the individual and the
actual world.\footnote[1]{S. \emph{Kierkegaard} has sought to maintain a
complete opposition between ethics and history, and from his standpoint
rightly. ``God could, without doing wrong and without denying his nature,
which is love, create a man endowed with capacities like no other, and
set him in an out-of-the-way place and say to him: Live through now, in
exertion like no other, the human; work, so that the half could be enough
to transform a contemporary age, but you and I, we remain alone about it;
all your striving shall have no significance whatever for any other human
being, and yet you shall, do you understand, you shall will the ethical,
and you shall, do you understand, you shall be enthusiastic, because this
is the highest.'' Uvidenskabelig Efterskrift p. 98. --- From the ascetic
point of view obedience, which is here the cardinal virtue, shines the
more brightly the smaller the result is. Whether one transforms a
contemporary age or waters a withered stick comes in any case to the same.
Cf. above p. 71.} This relation must not be made into a sham relation,
which it must always become under dualistic presuppositions. The miracle
conflicts with ethics just as much as with physics and logic.

There is a harmony between personality and the ethical tasks. Personality
is developed by work on the solution of the tasks, and conversely the task
consists, directly or indirectly, in the development of personal life. A
society consists in a plurality, a manifoldness, which is bound into unity
by peculiar laws. The more freely and independently the single members
move, the more different \opage{103}possibilities they realize, while at
the same time the unity is preserved and constantly gains a more intimate
character and higher validity, the more perfect is the society. The idea
of the human race as the kingdom of personalities, as a living unity of
individual forces, is therefore the highest ethical idea. It is to it that
the one who acts ethically, consciously or unconsciously, under the one
form or the other, looks when he weighs and appraises ends and means. It is
not the fruit of arbitrary speculation, but is an ideal that has formed
itself for human consciousness in the course of history\footnote[1]{Cf.
my essay: ``Forestillingen om Retfærdighed i dens historiske
Hovedformer.'' Nationaløkonomisk Tidsskrift. 2det Bind p. 81---116.},
thought fetching, in artistic fashion, the elements of the image of the
perfect from imperfect reality, an ideal whose partial, progressive
realization can be shown.

But progress does not go in a straight line, and not all parts of the
human race share in it. This might seem to be a weighty objection against
the whole idea of progress. Yet one must not, over the reactionary
periods, forget the essential tendency, the law, that reveals itself when
one compares a series of periods, especially as reaction itself proves to
be an important member in the development, which enriches it with new
experiences and firmer forms. All motion is rhythmic, but that does not
exclude its being progressive. ``Where many different
\opage{104}elements take part in the motion, the same state will never
quite return. A political reaction never brings back the old order of
things. The rationalism of our time is very different from that of the
last century.''\footnote[1]{\emph{Spencer}: First principles. Part II.
chap 10: The rhythm of motion.} And as regards the unequal participation
of the different races in progress, the case here is as within the animal
world, where the development of the type in some groups stops at a point
that for other groups stands only as a point of transition. The human
fetus shows at early stages a likeness to lower types. So it goes with
the civilization of peoples; the one people stops at a stage that is known
from other peoples' history as a standpoint overcome. Thus many oriental
societies have stopped at the stage of development where the civil and the
ecclesiastical government still coincide. One can perhaps even say, if one
disregards the highest stages of culture, that stagnation is the rule and
progress the exception.\footnote[2]{\emph{Maine}: Ancient Law p. 22.} Only
a small part of mankind has reached a higher civilization. But whereas the
more highly developed animal species is unable to give the lower species a
share in the results of its development, such a communication is possible
within the human world. European culture seems destined to become
universal, to take the other forms of culture up into itself, although the
meeting with these more primitive forms of culture cannot yet be said
\opage{105}on the whole to have had a very humane character. Europeans
still have toward the inhabitants of other parts of the world no small
degree of the self-esteem and national pride that the ancients cherished
toward the barbarians. Yet this barrier too will one day fall, all rivers
will unite in one great stream, and a single community of culture will
embrace all the inhabitants of the earth.

At the ethical standpoint one has no time to dwell on the thought of how
much has been accomplished (``wie wir's so herrlich weit gebracht''), just
as little as one has time to brood over what could have been otherwise.
The pressure of what remains to be done will even, as a rule, at any time
be so strong that men will be reluctant to admit that they are happier
than their forefathers. The good old time, of whose vices and sufferings
one has no living feeling, but in whose virtues and goods one believes, is
therefore as a rule extolled at the expense of the present time. Men have
for the most part felt themselves to be in the iron age and placed the
golden age in the past. Our time seems especially to be the time of
pessimism. This can be derived from several causes. The great political
and social reforms in the last century were greeted with expectations of a
golden age of freedom which have not been fulfilled --- among other things
for the simple reason that those who were emancipated, by the nature of the
case, still lacked the education that only freedom itself can give. The
great upswing in industry and commerce has not been accompanied by the
right distribution of all the goods that were thereby obtained. The
advanced \opage{106}intellectual culture has called forth a reflection and
a consciousness which the primitive peoples lack, and which can often
hinder the free and full enjoyment of what is valuable in life, while it
sharpens the sense for the disharmonious and for the shadow sides. In the
ethical respect, however, precisely hereby a great advance has been made,
in the direction of sympathy with the suffering of others. A personal
experience of pain is necessary in order that sympathy with the pain of
others may gain the right strength and extent. And the urgent
consciousness of what remains to be done presupposes the idea of the high
goal that was to be reached. Pessimism itself will thus, where it is more
than the fruit of the disappointed hopes of a refined craving for
enjoyment, lead to a clearer conception of the conditions of human life and
to the feeling of how important it is that work be done toward their
perfection. In the ethical principle, as it has been developed in the
foregoing, is contained the possibility of reconciling work and enjoyment,
the morality of duty and the morality of utility. The fundamental ethical
task is, as we have seen, to develop harmony within the individual
personality as well as between the single persons in society; but the
highest happiness consists precisely in such a harmony, in the development
of the personal powers in concord with one another.

--- The ethical relation will not disappear again as long as the
fundamental conditions of human life remain the same. It owes its origin
and its activity in mankind to definite conditions, \opage{107}but these
conditions are, as we have sought to show, the same as those on which the
existence and development of the race itself depend. The forms under which
the ethical appears may vary without its changing in its innermost nature.
Only if it coincided entirely with the relation to the external
authorities would its role be played out when emancipation triumphs; then
we should already now be witnesses of its process of dissolution.
Advancing culture leads men to look with irony on all authority, and if
with this they reject everything that grew up in its shelter in the course
of the ages, the ethical too will disappear. But in the blasé mood that so
easily arises in the transition from one great principle to another, one
must not see the final result of the course of development. A closer
investigation of the psychological and historical foundation of the ethical
leads to the recognition not only of its great and lasting significance,
but also of its inner connection with the fundamental laws of mankind.

% ===== FREEDOM OF THE WILL (pp. 108--146) =====
\chap{Freedom of the Will.}
\opage{108}The recognition of the ethical law differs from theoretical insight
and aesthetic contemplation in this, that it makes a demand on us. But here
the question then naturally arises how far man is able to satisfy this demand,
or the question of the freedom of the human will. This too is a problem that
concerns the foundation of ethics, for which reason we must still give a brief
indication of how it is to be understood from the standpoint set out in the
foregoing. It is with deliberate intent that this inquiry is undertaken only
now; for however disputed the questions with which we have occupied ourselves
in the preceding chapters may be, the problem of the freedom of the will
nevertheless seems to have a peculiar sting. This has its ground in part in
the fact that the psychology of the will is still so little cleared up; the
will is perhaps the most intricate side of man's nature, where very many
conscious and unconscious factors play a part, while \opage{109}the
development of cognition and of the life of feeling is simpler and easier to
follow. But it must surely be derived chiefly from the fact that people in
general believe they know extraordinarily well how matters stand with this
question, for which reason one is easily repelled by the divergent results of
a theoretical inquiry. The supposed knowledge hinders the growth of right
insight not only by excluding the need for it, but also by favoring feelings
that make it impossible to consider the matter in an unbiased way.

Wherever we seek to understand and explain a thing, we must go back to the
larger connection within which it comes into being. We must see it develop and
consider its relation to the other things that are links in the same chain.
This procedure we have followed in the preceding inquiries, and we will also
apply it in the present one. By thus seeing the thing develop, one more easily
discovers the elements contained in its nature; the different stages in the
development are marked precisely by new elements attaching themselves to those
previously given, and by the connection between the given elements taking on
another character. But herein lies, then, that in order to understand a thing
we must dissolve (analyze) it. As it stands in its finished and completed
unity it is incomprehensible. It must be broken apart, each of its parts must
be considered by itself, and we must then see whether we can understand the
working together of the parts, and so put together again what we have taken
apart. Research often fares as children fare with their toys, which they too
tear \opage{110}to pieces in order to study their hidden nature, but which
they are not able to put together again. But this danger is amply outweighed
when the process of dissolution (analysis) merely teaches us to know some new
elements; the first step is then taken and can draw more after it; we cannot,
after all, expect to reach the goal in one day. The immediate apprehension,
however, as a rule feels repelled and uneasy at this dissolution of what
stands for it as a completed whole resting in itself, and in every domain
inquiring thought has had to fight a hard battle to assert its right. It seems
as if all individuality were threatened by research in a struggle of life and
death. For the lifeless world one will nowadays no longer raise objection to
the application of analysis; one will not readily speak of individuals here.
Only to the imagination does a mountain or a distant island appear as an
individual being. But in the organism many still believe they have an absolute
individual, although \emph{Goethe's} words: ``Kein Lebendiges ist eins, immer
ist es vieles'', have received complete confirmation through the results of
science, and hold not only for the organism as a whole, but also for its
parts; even the cell, in which for a time one believed one had found the real
individual, is a little world, and so on to infinity.

Must we then wholly give up the concept of individuality? Does not nature then
dissolve for us into a chaos, where innumerable elements float about among one
another and are joined only for a moment to form apparently fixed and
completed shapes? \opage{111}This conclusion springs only from the impatience
of the imagination at having to give up the indissoluble unity with which it
has endowed things. With far greater right we must turn the relation around
and say that the whole activity of nature is individualizing. % misprint in the print: "Virkomhed" for Virksomhed (activity)
The analyzing inquiry seeks only the elements that in nature are joined into
peculiar wholes. At all the different stages of the life of nature there form
centers or nodal points, from which particular activities go forth. The nature
of each of these peculiar activities depends on the reciprocal relation
between the given center and the surrounding conditions. Already within the
lifeless world there form smaller circles, which modify in a peculiar way the
influences coming from without. Think, e.g., of the state of aggregation of
substances and its dependence on temperature and pressure; in this definite
relation to the conditions lies something individual. But greater
independence, and therefore also the possibility of a further-reaching and
deeper interaction with the surrounding world, we find first in living beings,
which therefore mark a higher stage of individualization. The living being is
able, through the freer and more intimate relation between its elements, to
adapt itself in far higher degree to its surroundings --- indeed, at the
higher stages, even to adapt the surroundings to itself. Recent science, which
lays a special weight on the relation between the organism and its conditions
of existence, has thereby paved the way for a fuller conception of
individuality, which was formerly often regarded as an \opage{112}isolated
unity. It is in the manner in which there is reaction against certain
conditions that individuality shows itself. The individual is not substance,
but function. The greater or lesser richness and manifoldness of the
interaction furnishes the standard for how high a being stands in the series
of things. And it is through this interaction, too, that the transition,
according to the most probable hypothesis, takes place from the lower to the
higher forms. The struggle for existence consists in the individual's
asserting itself in its peculiarity under the given conditions, and is thus a
struggle for individuality.

From this point of view the mental life, too, stands in close connection with
the whole development of nature. It marks the highest stage of
individualization that we know. Consciousness is the condition for the
reciprocal relation to the world taking on a comprehensive and deep-going
character. The different influences from without must be gathered, remembered
and combined; distinctions must be made between different kinds of influences,
in order that the reaction may correspond exactly to that which calls it
forth. But that gathering and this distinguishing can take place only in a
consciousness. The different stages of the development of conscious life ---
``from the intestinal worm in an organic tissue to a Newton or Shakespeare,
whose thought embraces the world'' --- are also stages of the development of
individuality.

The human individual has, through its whole organization, the possibility of
entering into a richer reciprocal relation with the surrounding world than the
other living beings. Yet originally it, too, is only a side track,
\opage{113}along which the movement from without is conducted further, without
anything happening in the individual itself other than a mechanical
propagation of the communicated impression from the sense organ to the central
organ and from this to the organ of movement. This is the case with the
so-called reflex movements. Here movement follows immediately upon impression,
as when the pupil contracts under the impression of strong light, or when
someone, without himself knowing it, chases away a fly that has settled on his
cheek while he is absorbed in his work. Such movements, which are direct
discharges or reflexes of the influence from without, can therefore also be
carried out when the individual has been deprived of the most important organs
of conscious life, as the many generally known experiments with frogs and
salamanders prove. Robin even saw the arm of a beheaded criminal begin a
warding-off movement when he scraped him on the chest with a scalpel. Even
where the movement proceeds from the individual itself, as in the spontaneous
and purposeless movements in which the nervous system seeks an outlet for its
increased energy (see above p. 11), it still has a purely mechanical
character. And yet in the capacity for spontaneous and reflex movement we must
see the germ from which the will develops, without its being possible to
indicate the definite point where the involuntary ends and the voluntary
begins. Of will we speak only when there has formed a fixed and coherent
circle of ideas and feelings by which the individual's action is determined,
so that the impression received or the spontaneous impulse is not immediately
converted into outward-going \opage{114}activity. Before such a fixed core has
formed in consciousness, we do not ascribe to the individual any real self
either. But the point at which such a self can be said to be present is
reached through a long development, which begins with very few and very
obscure ideas and feelings. Some physiologists have even assumed that a
sensation takes place even in the pure reflex movements, although this does
not come into our consciousness. So much is in any case certain, that no sharp
boundary can be drawn between the involuntary and the voluntary movements.
Even if sensations or ideas appear between the impression and the action that
follows from it, the action can still be called purely mechanical in so far as
the real self, i.e. the whole circle of thoughts and feelings that otherwise
rules in the individual's consciousness, does not come into play at all. This
is the case with all actions performed from habit, where we often act
according to a quite distinct idea without this idea entering into interaction
with the rest of the circle of ideas. It is on the whole the case in
``unguarded moments'', where attention is not awakened; we then do something
other than we really want. But the greater the significance the interval
between prompting and action acquires, the more reason there is to speak of a
will. In this interval the prompting can be brought into connection with the
rest of the content of consciousness and act upon the mental state as a whole.
What thoughts and feelings \opage{115}it thereby awakens in the individual
case depends partly on the general psychological laws of the association and
reproduction of ideas, partly on the individual's inborn temperament and his
character as developed in the course of life. An effect can never be derived
from a single cause, but always presupposes a whole series of conditions. The
action becomes a consequence of the manner in which the individual's nature is
modified by the impression or the prompting. The more restricted the circle of
ideas is, the fewer possibilities present themselves to consciousness, and the
more immediately the impression leads to action. The harder it is to set the
circle of ideas in motion, so that images of memory can emerge, or thought
turn toward the future, or in general toward objects that lie far from the
situation given at the moment, the less developed is the will. Children and
savages are therefore children of the moment, and let themselves be carried
away immediately by the first prompting that comes along. Travelers emphasize
as the most characteristic trait of savages their lack of capacity for steady
and persevering work whose reward only the future is to bring. Their feelings
are quickly awakened, but last only a short time. They have, as has been said,
the understanding of children and the passions of men. Without memory there is
no coherence in conscious life, and therefore no real self either. Instead of
distinguishing between voluntary and involuntary movements, one could
therefore distinguish between movements that presuppose memory and reflection,
and such as are performed \opage{116}immediately and without forethought.
\emph{Wundt} has rightly emphasized that animals which have been deprived of
the central organs of consciousness lack memory, which is seen when one lets
some time pass between the experiments; the same unsuitable movements are then
made anew before the suitable ones are found. But the boundary remains always
movable, since actions that were first performed with reflection are gradually
performed wholly without forethought.

The will thus presupposes deliberation. Prior to the willed action there goes
on a struggle in consciousness between different possibilities, different
promptings to action, and the object of the will becomes the prompting that
triumphs over the others. The richer and more strongly developed conscious
life is, the more the action, too, can be called the individual's own,
according as it is decidedly determined by his nature. The fixed core that has
formed in consciousness will assert itself in every essential action that is
to be performed, and stamp it with its impress. In other words, a character, a
personality, has formed. Reason has, as Hamlet says, mingled with the impulse
of the blood, ``so that Fortune cannot, as on a flute, play on him whatever
melody she pleases''. The decision is not an immediate discharge, but the
fruit of a process that has taken place in the individual's innermost nature,
and to which all the elements in it have made their contribution.

Every decision proceeds from the \opage{117}existing psychological
presuppositions with the same necessity with which effect everywhere follows
cause. But it is not any external power that overwhelms man; his own nature,
his own self, asserts itself in the strongest motive in a particular form
determined by the circumstances. It is man himself who acts. The agent may
have a feeling that none of the motives that have asserted themselves to his
consciousness contains the sufficient cause of the action. But that only shows
that one must go deeper to the bottom to find the cause of the action. Behind
the conscious motives lie the unconscious, inherited or acquired
predispositions, which exert a decisive influence on our action. Hence it
comes that we learn to know our own natural character only gradually, through
experience; by what we do we learn, as \emph{Schopenhauer} says, what we are,
and we are often surprised in a high degree by the discoveries we make in this
way.

In man's nature there lies, as we have seen earlier, an original impulse to
activity; how it expresses itself is determined by the organization received
from the hand of the race and of nature. The individual's own experiences and
actions also leave behind them traces and after-effects, which assert
themselves as predispositions and dispositions in certain definite directions.
We thus never begin absolutely from the beginning in our action. As links in a
great connection and in an infinite series of development we have our definite
point of departure and our definite \opage{118}presuppositions, which are one
with our individual nature, our original character. Therefore the freedom of
the will is wholly illusory if by it one understands a capacity to begin an
action absolutely, to be cause without being effect. A cause that is not an
effect, that is thus its own cause (causa sui), must be an absolute being,
raised above all conditions, torn loose from the lawful order of things. But
such a being is incomprehensible to us. Acute defenders of the absolute
freedom of the will have also seen this; for them it is a postulate that
cannot be grounded. But however much one recognizes the limitation of human
knowledge, one cannot thus throw down the bridge in the midst of the world of
phenomena. If the connection here is burst at a single point, it will be all
over with all lawfulness and all knowledge of nature. Either one must make man
a god, or else one must recognize that he arises and develops according to
definite laws, which to a certain degree must be or become accessible to us.
Experience now shows us clearly enough that he is a finite being, even if the
individualizing process of nature has reached in him what we must regard as
its highest point. Individuality rests here, as everywhere, on there forming a
smaller, relatively closed circle within the larger, and on the influences
coming from without being modified by the laws that hold for this smaller
circle.

But just as surely as we first learn to know our character from our actions,
just so surely is it \opage{119}also the case that our actions conversely
contribute to determining and developing our character. It is a physiological
fact that the structure of an organ is developed by exercise. The result of
activity is not merely the work accomplished outwardly; it is also the
increased skill, the more developed disposition in a certain respect, that the
individual has thereby attained. Every feeling, every decision, every action
adds a stone to the building of man's character. And when such a spiritual
development has now been reached that the individual can become conscious of
this, as he begins to fulfill the Socratic demand: Know thyself! --- then he
will also be able to seek to guide the course of his ideas, feelings and
decisions in the direction he regards as the best. Man stands over against
inner nature just as he stands over against outer nature. He cannot work
miracles. He cannot intervene in the course of things to guide it according to
his purposes, except in so far as he can learn to know the conditions under
which things arise and act: if he can bring about these conditions, then he
can also guide the phenomena of nature according to his will, e.g. divert the
lightning bolt, or he can at least foresee their occurrence and take
precautions against their consequences. It goes similarly with him in regard
to his own nature. Here, too, he cannot create or work miracles. But if he has
sufficient knowledge of himself, that is to say, of the manner in which his
nature is affected by circumstances and \opage{120}conditions, then he can
seek to change these conditions in order thereby also to change his own ideas
and feelings and guide them in a certain direction. How far one can get along
this path only experience can show. The limit will always be staked out by the
original predispositions in man's nature; and only fantasticality will venture
beyond it. Here, too, mastery shows itself in limitation. But much that, seen
immediately, seems impossible can nevertheless be reached, if not directly
then indirectly. We cannot act immediately on the movement of the heart, but
by holding our breath and compressing the chest, or by chemical influence, we
can modify its movement when we will. --- ``If one cannot leap over the
Atlantic Ocean as over a ditch, one can still sail over it; and although
Columbus was assured that it could not be done, and that anyone who attempted
it must necessarily perish, he nevertheless did it and did not perish. In like
manner transformations of character are possible, but, to be sure, only within
the limits of a certain given natural endowment, and here, too, the execution
requires time, courage and perseverance''.\footnote[1]{\emph{Fortlage}:
Psykologiske Foredrag. Kbhvn. 1876. p. 187.}

Everything depends on attention and interest being awakened. In the domain of
cognition man is at first guided immediately by the impressions and by the
associations of ideas that they call forth, but these often have nothing
whatever to do with the \opage{121}actual connection of things. When he is led
by experience to see this, attention is awakened, and he strives to form and
connect his ideas in such a way that they correspond to reality. Consciousness
is nowhere wholly passive. Activity is the condition for receptivity.
Attention, which is of such great significance for cognition, is already a
will.\footnote[1]{``Apperception [the mind's attentive grasp of something] and
the impulse to voluntary movement are only different forms of the awakening of
the will. --- --- The voluntary movements most probably proceed from the
foremost parts of the brain, while the central organs of sensation seem
chiefly to be found in the parts of the surface of the brain that lie further
back. On the other hand, one can hardly doubt that the higher mental functions
are bound chiefly to the development of the forebrain. This connection first
becomes intelligible when we consider that the points of origin of the
innervation of the will at the same time govern the centers of sensation, and
thus determine not only the movements but the apprehension of sense
impressions and the course of the reproduced ideas.'' \emph{Wundt}:
Physiologische Psychologie. p. 765 f.} Painful experience is often the
condition for its being awakened. In the domain of action as in that of
cognition, man becomes wise through harm. The dominion of the immediate
natural impulse must often lead to decided conflict with what man holds fast
as his right task, his true end, before a penetrating consciousness of the
significance of the action can arise. When thought has awakened and the
consciousness of a higher law has developed, a new impulse has awakened in
man, which strives gradually to master \opage{122}his action. But it is only
step by step that the opposition between the two tendencies becomes clear;
temptation itself will then, according to the law of the association of ideas,
awaken the idea of the law against which it strives, and thereby make a
struggle possible. The passion awakened by the impression of the moment will
long continue to obscure the consciousness of the essential law.

``Noch ist, bei tiefer Neigung für das Wahre,\\
Ihm Irrthum eine Leidenschaft''. (Goethe).

A condition for the idea of an essential law being able to form is, as we have
seen earlier, that the individual has lived its way into the thought of the
larger whole to which it belongs. This happens through education, and
therefore education is also the presupposition of self-education. Man must be
guided in order to be able to guide himself. Freedom in the sense of a
capacity for self-education can come into being only where this condition is
fulfilled. Here again we see how closely the single individual is bound to the
whole; what he is capable of depends on his relation to it. It is only the
fully developed individual that can be free in the sense in which we here take
this word.

In defense of the assumption of absolute freedom, of the capacity to be cause
without being effect, one appeals as a rule partly to the general
consciousness, the immediate feeling, partly to remorse and to accountability.

The immediate consciousness and feeling can \opage{123}decide nothing here. It
teaches us only what we feel and imagine, but not how the feeling and the idea
have arisen, or in what relation they stand to the actual connection of
things. It grounds no theory, neither in the one direction nor in the other;
it is only reflective thinking that poses the problem and seeks an
explanation. The immediate consciousness is not wrong in asserting that we
ourselves are the cause of our actions; what it is wrong about is that we
should be absolute cause, cause without being effect. It assumes this because
it does not feel itself compelled or set in motion from without when it
decides, in the same way as it assumes that the earth stands still because the
motion is not noticed. We do not ourselves notice the motion of our will, any
more than we notice that of the earth; in both cases we infer the motion.
Nowhere is there anything that moves without itself being moved.

We have shown earlier how the higher stage of life passes judgment on the
lower. Just as the history of the world is the judgment of the world, so the
history of the individual man is the judgment on himself. We measure our
earlier action by the more perfect knowledge and feeling of the right that we
have since attained. In the calm and harmonious mood we pass judgment on what
we did while passion ruled us. Thus we apply an ideal standard to ourselves, a
standard taken from another standpoint than the one from which the action was
done. The pain that arises from the feeling \opage{124}of the opposition
between our actual action and the ideal action we call remorse. Remorse thus
becomes intelligible and is seen in its great ethical significance even if we
assume no absolute freedom. It consists, on our view, in the penetrating and
painful feeling of \emph{not yet} having got further. How strong and intense
this feeling is depends on the strength and clarity with which the idea of the
ideal, of the obligating law, appears. Remorse points forward; it is the birth
pangs of the ethical character, a testimony to our predisposition to
development. If it were to presuppose that we could also have acted otherwise
under the given circumstances, it would in any case be a wholly idle
consideration to dwell on, and \emph{Fichte} would be right in his energetic
remark that we have no time to repent. The value of remorse rests partly on
the idea of the right becoming clearer and stronger through the deep feeling
of the wrong, partly on the closer knowledge of himself that man gains through
it. What has happened cannot be undone, but the new possibilities that present
themselves through self-knowledge act as a spur and drive forward. From the
standpoint of humane ethics it may therefore hold: felix culpa! (happy guilt!)

As for accountability, finally, it rests on man's being so developed that
there has formed a self, a fixed core in his consciousness, which can
determine his action. On the other hand, it presupposes no absolute freedom.
However \opage{125}much I may be able to give an exhaustive explanation of my
action, may see the connection and necessity of all the motives, I must
nevertheless reject it and impute it to myself when I see its disagreement
with what I recognize as the highest law of human action, and when I was in
full possession of my reason when I performed it. The motives were, after all,
not foreign to me; they were precisely the expression of my nature and
proceeded from it with necessity under the given circumstances. How else could
the action be my own? If at a single moment I could act in such a way that the
action had nothing whatever to do with my actual nature, did not proceed from
it with necessity, my personality would split into many parts. It is true that
one rightly distinguishes between man's real self or true will and the manner
in which he shows himself in his actual action, and the Apostle Paul indeed
even says: ``I do not do the good that I will, but the evil that I do not
will; but when I do what I do not will, then it is no longer I who do it, but
the sin that dwells in me.'' (Epistle to the Romans VII, 19---20). But one
must not confuse a figurative opposition with a real separation. What we call
the real or true self could also be called the ideal self, the character as it
would be if it were freed from its imperfections. But it has no existence
except in the actual self itself. Our faults are ``the faults of our
virtues'', and it is \opage{126}only to excuse ourselves that we take our
``better self'' for our real self.

Accountability, like remorse, points forward, not back. That I recognize
myself as accountable with regard to a certain action means that I feel the
obligation to work further in order to get beyond the imperfection that the
action has brought to light. I feel myself liable to renewed education.
Freedom in the sense of a capacity for self-education is, as we have remarked,
possible only when an education has gone before. It is a capacity that
presupposes certain definite conditions. When the clarity and liveliness of
consciousness, whether on account of youth, of insanity, or of passing
abnormal states, are not present in such a degree that man can ``look forward
and back'', consider the nature and consequences of the action, then the
juridical law does not regard him as accountable either; and it is a sign of
the growing psychological insight that the newer penal codes lay greater
weight on the influence of such circumstances than the older ones (compare
e.g. the Penal Code of 1866 with Christian V's Danish Law). A normal and fully
developed conscious life is the condition of accountability. This condition is
not present in children, whose education is not completed, nor in the insane,
or rather, it is found in them only at a lower stage. For the child, too, is
called to account and punished for what lies within its ethical domain; and as
for the insane, \emph{Maudsley} has aptly
% continues on p. 127
% continues from p. 126
\opage{127}remarked that one must assume a certain responsibility in them,
since in the asylums one acts upon the insane through the usual motives,
although one does not punish them as fully responsible when these motives do
not gain power over them.\footnote[1]{\emph{Maudsley}: Le crime et la folie.
p. 190. --- Cf. \emph{Griesinger}: Pathologie und Therapie der psychischen
Krankheiten. 2te Aufl. p. 479: ``In the asylums the patient's outward
restlessness and the audible expression of his morbid impulses are checked
already by the example of the other patients, by the prevailing spirit of
peace and order; he is drawn of himself into the calm movement of the whole
house; possible resistance on his part is met not so much by direct
compulsion as by his own feeling of submission to the imposing power of the
whole. Here he finds forbearance and attention, but also notices that
obstinacy would not help; he soon learns to comply with the physician's
arrangements and sees how the manner in which he is treated, the measure of
freedom and enjoyment that can fall to his share, depends on his composure
and conduct. Thus he finds here essential help toward self-control.''}
Accountability is thus relative, varies with conditions and circumstances.
Fully developed (``grown-up and sound of mind'') persons are compelled by
society to obey its laws. This too is an education toward a definite end and
presupposes accountability only in the sense that there must not be anything
in the persons that excludes the possibility of reaching this goal. He who,
under otherwise normal conditions, proves in deed to lack respect for the
law of society must be brought to recognize it by feeling its power. Our
psychological-ethical standpoint thus leads us to agreement with
\opage{128}the manner in which Professor Goos in his ``Indledning til den
danske Strafferet'' (p. 60---64) has grounded the penal authority of the
state.\footnote[1]{The present treatise was essentially finished when Prof.
Goos's book appeared, but it was a great encouragement and confirmation to
me to find a kindred line of thought in this acute work.} From a juridical
standpoint one certainly cannot regard education in the sense of ethical
improvement as the end of punishment, since the power of the state has no
right to intervene in the individual's innermost personal life. Nor, after
all, does the medical treatment of the insane aim at improving them, but
only at restoring the old self in its health; and \emph{Griesinger} finds in
this the difference between education and the treatment of the insane. What
society must demand is practical recognition of its authority, and by
punishment such recognition is compelled where it has proved to be lacking.
By punishment, as Prof. Goos says, the criminal's character is bent to
obedience to the law. ``The encroachment on the person's legal goods, the
suffering that punishment contains, aims at supplying the person with new
and stronger motives for the disciplining of his character to obedience to
the law, after it has proved that the disciplining left to himself has not
brought about what the power of society requires in this respect.'' Prof.
Goos draws from this result the conclusion that, as a presupposition for the
justification of punishment, accountability is not required in the sense of
a capacity for absolute freedom, ``freedom of the will in the sense that the
application of the law of causality \opage{129}to human acts of will, or the
conception of the relation between motives and decision as the relation
between cause and effect, is excluded.'' The necessary presupposition is
only that disciplining by human effort is not impossible, and it is not so
as long as it is still possible to awaken the man's self-feeling, his
attention to himself, his capacity for energy and resignation. Only he can
be liable to education in whom it is possible at all. The crime gives
society the right to educate the criminal. The criminal's relation to
authority sinks back, on account of his own deed, to the starting point;
since he does not feel reverence, he must feel fear.

We have earlier insisted that the single individual can never be isolated
from the larger context in which it arises and develops. Nevertheless it was
such an isolation that was presupposed by the manner in which punishment was
regarded in older times. Whether one wished by it to threaten and deter, to
requite eye for eye and tooth for tooth, or to propitiate an angry deity,
the criminal stood as the enemy of the whole society, as one toward whom one
no longer had any duty. In great part this way of regarding it doubtless
sprang from the passionate feeling of revenge, which sees the offense as
absolute and therefore admits no limitation; but the dogma of the absolute
freedom of the will has also exercised its influence here. It dissolves the
race into individuals, into absolutely independent starting
\opage{130}points. As soon, on the other hand, as one sees the deep
connection between the individual, in the whole range of his nature and his
character, and the whole race and society, it must also appear impossible
that anything whatever should be able wholly to loosen this connection. As
the crime itself often points back to disproportions and defects in the life
of society, so on the other side it must be the task of society not to
annihilate, but to help forward and develop the germs of truly human life
that may be hidden even in the apparently most hardened mind; it thereby
works for its own cause, its own development.

The freedom of the will, as we recognize it, does not conflict with
necessity. The validity of the natural laws embraces ethical action too. By
the freedom of the will we understand the state in which man's reason and
feeling, and his spiritual capacities in general, are so developed that a
single impression, a single feeling or passion, does not completely dominate
him or prevent reflection and deliberation, the application of the standard
that stands for him as the highest. Freedom is thus here the opposite not of
necessity but of blindness and chance. It makes self-control and
self-education possible. But whether we regard freedom as a capacity or as a
state, it is always relative. Our freedom consists in a striving to become
free. Few are the points at which the individual really unfolds full
activity. The further his development in this respect has
\opage{131}advanced, the greater too will be the contribution that his
self-conscious action makes to the formation of character. And when the
feeling for the significance of this contribution is sharpened, when
attention is drawn with force in this direction, consciousness forms for
itself an ideal model in which the contribution from the conscious action of
the personality becomes so great that the other elements vanish in relation
to it, --- the ideal of perfect freedom, of what the human personality would
be if all hindering conditions and influences were thought away. This ideal
is only a special form of the general ethical law, which, as shown earlier,
consists in the idea of perfect human life as that which develops through
the ages. In this idea the image of action springing freely, i.e. from man's
conscious personality, is also a necessary element; it expresses the form of
activity that accords with man's place in the series of natural beings.

But this ethical \emph{ideal} the common conception confuses with the
\emph{actual} will, which is always conditioned. Perfect freedom we conceive
by abstraction and idealization. The thought of it acts, like every ideal,
as a spur and incitement; but it rests on an illusion when one transforms it
into a capacity that is supposed to be present in all and under all
circumstances. Such an illusion is nothing unique, but has in every domain
hindered the conception of natural development. One refuses to recognize the
\opage{132}perfect as the result of a development and regards it as being
there from the first. Thus one conceives the soul as a being by itself that
is implanted in the body, and the fully developed body as already present in
the germ en miniature. What development itself first leads to must in one
way or another have been there from the beginning. One has thereby come to
ascribe to all men at every moment, as a capacity, what is in itself only
the goal of a long development.

Most attempts to maintain the absolute freedom of the will do not exactly
have such a scientific form that they invite closer discussion. Only one
form of the absolute doctrine of freedom will it be worth while to bring
forward here, namely \emph{Kant's}, since he is the most eminent of the
inquirers to whom one will be able to appeal in this matter. Here there is
now, in the first place, the remarkable circumstance that Kant himself with
great emphasis stresses the natural necessity of human willing and acting,
in so far as man's being is accessible to us through experience. In
psychology Kant is decidedly a determinist. ``All the actions of man in the
world of experience,'' he says in ``Kritik der reinen Vernunft,'' ``are
determined by his empirical character and the other co-operating causes
according to the order of nature, and if we could investigate all his
voluntary actions to their ultimate ground, there would not be a single
human action that we could not predict and recognize as a necessary
consequence of the preceding conditions. With \opage{133}regard to this
empirical character there is therefore no freedom, and yet only according to
this character can we regard man when we merely wish to \emph{observe}, and,
as happens in anthropology, investigate physiologically the moving causes of
his action.'' It is the same fundamental conception that we have sought to
carry through in the foregoing. Kant was so permeated by the spirit of
determinism that in his doctrine of method he maintains that even the
wildest hypotheses, if only they keep within the domain of natural causes,
are preferable to the appeal to the supernatural, which serves ``lazy
reason'' only as a pretext for giving up the labor of finding the causes by
way of experience. But according to Kant's whole conception of the ethical,
this does not belong at all to the world of phenomena, of experience. The
ethical law, which makes itself known in our practical reason, can, as we
have seen earlier, receive no grounding at all according to Kant's doctrine.
It speaks to us from a higher world than that of phenomena, absolutely
independent as it is of all experience. In so far as man's will is
determined by the ethical law, he is thus raised above the world of
experience. Besides his empirical character, which can become an object of
experience and observation, man therefore also has an ``intelligible''
character, according to which his action is guided by reason. ``Reason
itself,'' says Kant, ``is no phenomenon and not subject to any of the
conditions of sensible experience.'' Through it man gains the capacity to
act with absolute \opage{134}freedom, i.e. independently of all motives and
promptings given in experience.

One will see at once that Kant by no means assumes any decision or willing
without a necessary cause. What he understands by freedom is a determination
by ideal grounds (reason, duty) as opposed to determination by empirical
grounds (the impulses, pleasure and pain). As regards the ``intelligible''
character too there is a necessary relation between motives and actions, for
which reason Kant in his ``Kritik der practischen Vernunft'' (§ 6) sets
himself the following problem: ``Supposing that a will is free (sic), to
find the law that alone is fit to determine it with necessity (sic).'' He
further investigates carefully ``how the ethical law can become the
incentive of the will,'' and finds that this happens through the feeling of
respect and reverence. We have here, then, only another or higher kind of
necessity, not freedom in the sense in which we have combated it above. We
too have conceived freedom as an ideal that is not immediately given. The
difference between our doctrine and Kant's consists in this, that he sees in
freedom determined by reason something that lies beyond the world of
experience, whereas the freedom of which we speak stands precisely as the
goal of man's striving and development in the world of experience. We can
console ourselves here with the thought that we show greater fidelity to
Kant's philosophical principle than he does himself. His greatness rests on
his having demonstrated the subjectivity of our knowledge and its limitation
to experience. Only \opage{135}in the practical law of reason did he think
he had an exception. But by what right, then, does he call this law a
``fact''? A fact must surely be accessible to experience. When he says that
``reason itself is no phenomenon,'' how then does he come to know it, since
experience, according to his own principle, is the only source of knowledge?
And when he says that ``reason is not subject to the conditions of sensible
experience,'' how then can reason come into being and
develop?\footnote[1]{Kant here comes into contradiction with himself, since
elsewhere (Kritik der pract. Vern. § 7 Anm. 2) he says that practical reason
is ``a naturally acquired capacity.''} We know no other reason than that
which is the result of man's spiritual development in the world of
experience. It is thus itself a phenomenon and must be conceived according
to the laws valid for all experiences. The reason why Kant placed the
ethical law of reason and freedom in another world than that of experience
was their ideal character in relation to the other motives, capacities and
powers in man. It is a remnant of dogmatism in Kant that he here considers
it necessary to think of the ideal as existing. In reality, that man,
ideally regarded, is free means nothing else than that it is his task and
end to acquire freedom. What could it help that the ideal existed in a world
beyond, if it could not be striven for and realized in the world we live in,
the world of experience?

\opage{136}If one should object to this that such an ideal cannot possibly
be maintained, since experience constantly demonstrates our lack of freedom,
it must be remarked that any ideal whatever by its nature points beyond
reality. Without this opposition it would never come to action. The impulse
to act arises through the idea of a state that is not present but whose
image memory preserves; when conscious life has reached a certain
development, the idea of a more perfect state takes the place of actual
memory, and this is then, to be sure, often confused with actual memory;
hence the myth of ideal innocence and happiness in a vanished paradise.
Every impulse points beyond the given, and the ideal is the completed and
harmonious expression of what already dawns partly and obscurely in the
primitive manifestation of impulse; therefore it has its root in reality and
yet leads beyond the reality given at any time. Only by building on the
ground of determinism will ethics be able at once to maintain its ideal
standpoint and to preserve its close connection with the real. Therefore
great ethical thinkers too have sought to carry through the deterministic
line of thought. Reference may be made here in particular to the English
school (\emph{Hume}, \emph{Mill}, \emph{Spencer}), \emph{Herbart},
\emph{Schleiermacher}, \emph{Spinoza}, the \emph{Stoics} and \emph{Kant},
when we take him according to the spirit, not according to the letter.

One must go back to the founder of humane ethics, \emph{Socrates}, and the
movement proceeding from him, to find the explanation of the
\opage{137}difficulties in which the question of the freedom of the will has
become entangled. Socrates led human action back to self-knowledge and
maintained knowledge to be the essence of virtue. Passion became for him
ignorance or incorrect knowledge. From this sprang the definition of the
will which since \emph{Plato} and \emph{Aristotle} has been the prevailing
one, while through Christian theology it also permeated the popular
conception: the will is impulse guided by reason. Among the Greeks, in
accordance with their whole naturalistic and aesthetic view of the world,
the original endowment, the temperament and natural impulse were indeed
always stressed as essentially intervening in and conditioning man's action.
Plato derives all wickedness from a morbid bodily constitution and bad
education and has a gaze, surprisingly sharp for his time, for the nature of
insanity and its influence on the will. Aristotle holds fast to a happy
natural endowment as a necessary presupposition for ethical development. But
the Christian view of man in its severity took no account of such
conditions. It regarded responsibility as man's innermost being. Even the
obscure impulses and stirrings that dominate man before conscious life has
developed to clarity and strength are regarded as sprung from a will that
can be called to account. ``All emotions,'' says \emph{Augustine}, ``are
nothing but wills. For what are impulse and pleasure but will, in that we
approve what we will? And what are fear and sorrow but will, in that
\opage{138}we abhor what we do not will?''\footnote[1]{De civitate dei. XIV,
6.} ``By the sensual, the carnal,'' says S.
\emph{Kierkegaard},\footnote[2]{Kjærlighedens Gjerninger (3rd ed.) I, 60.}
``Christianity understands the selfish; nor can any conflict be thought
between spirit and flesh unless there is a rebellious spirit on the side of
the flesh, with which the spirit then strives; thus no conflict can be
thought between spirit and a stone, or between spirit and a tree.
Consequently the self-loving is the sensual.'' For the consistent Christian
view there is hidden a will in every mental stirring in man, just as behind
every natural event there lies a divine act of will. Here, then, two
absolute freedoms come to stand over against each other, of which the one is
the almighty, all-creating, over against which the other seems bound to
become nothing. This difficulty, which the Apostle Paul strikes to the
ground by holding up to man his absolute impotence (whereby he thus
precisely concedes the difficulty), Augustine seeks in his acute way to
solve by placing the will outside nature: ``God is the author of all nature
and all powers, but not of all wills. Evil has no nature; the evil will does
not spring from any cause whatever: it is, after all, only through the will
itself that evil comes into being!'' If the will is not nature, it also lies
beyond all corporeality. His chief objection to Plato's doctrine of the
influence of bodily \opage{139}constitution on the will Augustine takes from
the fact that the author of sin has no body. The ultimate source of evil is
a will without nature, which chooses evil with full knowledge and
consciousness.\footnote[1]{De civ. dei. V, 9; XI, 9; XII, 6; XIV, 3.}

To such consequences one is led by holding fast to a conception of the will
whereby an ideal state is immediately presupposed in man as an operative
capacity at every stage of his development and in every situation in which
he finds himself. We end here in the psychologically self-contradictory: for
the full knowledge and consciousness of the nature of evil, which is a
presupposition of the diabolical standpoint, would precisely make a willing
of evil impossible. It is impossible that man, in one and the same spiritual
state, should be able to choose the one just as well as the other of two
opposite things. Pride, which is supposed to be the beginning of sin, and
which consists in the self's choosing to be its own principle instead of
having its principle in a higher power, has, just as much as any other trait
of character, its natural origin, and the actions to which it leads spring
forth according to the general psychological laws. The theological
conception stops at the fact that the self wills to make itself the first
and to withdraw from the relation to the higher power. The psychological
conception goes further back and investigates why and how the self comes to
will this. \opage{140}The value of a poetic portrayal of a character rests
on how far it gives us the full explanation of all decisions and actions.
Even a character such as \emph{Shakspeare's} Richard the Third, in which the
poet has perhaps overstepped the limit of psychological truth and
probability, finds its reconciling explanation in the very first speech.
Contempt for the effeminacy and wretchedness around him, together with the
feeling of his own bodily and mental ugliness, which makes it impossible for
him to take part in the prevailing frivolous joy of life, while at the same
time he feels within himself an energy that could overcome all barriers, and
by which he stands high above those who regard him as an outcast, --- all
this begets the defiant resolve to play the devil's part:

``And therefore, since I cannot prove a lover,\\
To entertain these fair well-spoken days,\\
I am determined to prove a villain\\
And hate the idle pleasures of these days.''\\
% Shakespeare's English (Richard III, I, 1, 28--31); Høffding quotes the Danish:
% „og derfor, da jeg ikke kan som Bejler / behage denne tungerappe Tid, /
% saa har jeg sat mig for at spille Skurk / og hade disse Dages tomme Glæder.“
The poet has moreover clearly seen that in the psychological grounding one
cannot stop at the single individual, but must go beyond it to the whole age
and the preceding generations. Richard stands as the product of the
struggles of the red and the white rose, of all the cruelty and treachery
they begot. The passions have their natural history, and the ethical
consideration cannot leap over all the intermediate links by which character
and action first become intelligible, --- intermediate links which are
therefore properly parts of the very essence of \opage{141}character and
action. A French author has said: tout comprendre, ce serait tout pardonner,
and rightly, in so far as the action, by being seen in its necessity, no
longer stands as the purely arbitrary. But although the ethical
consideration cannot conflict with the psychological, it nevertheless does
not coincide with it. In psychology we look back, seek the preceding causes,
derive the later mental states from the earlier. In ethics we look forward,
measure moods, decisions and actions against a certain model, and ask only
about their harmonious or disharmonious relation to it. If it turns out that
they do not accord with the norm, then the task becomes to find the cause of
this discrepancy, and, when it is found, to investigate how far it is
possible to clear this obstacle away.

\emph{Hobbes} has aptly remarked that if the Aristotelian-scholastic
definition of the will as rational impulse (appetitus rationalis) were
correct, there could be no voluntary action that was contrary to reason.
What this definition overlooks is that impulse as well as reason must run
through a long development before they can unite and form rational will. One
cannot therefore regard this as a capacity that belongs to everyone from the
first, so that it depends on himself whether he will use it. In a certain
sense man never acts against reason, namely against the reason he has at his
disposal at the moment of action; but he can act against the reason he
\opage{142}ought to have, and which it is the goal of his striving to make
his own. From the fact that one cannot, it does not follow that one shall
not. Evil consists in man's being \emph{not yet} good, not in his having
ceased to be so. The degree of character and development of reason that man
has attained can be obscured and diminished when the right motives and
impulses are not maintained and when attention is dulled. But precisely this
surely demonstrates that he has not yet attained such strength, coherence
and consistency in his character that he can hold fast what he has acquired.
We know, as already remarked, our character only from experience, and it is
never concluded. Herein lies, on the other hand, also this, that we can
never know with certainty whether we have used all our powers. As we are in
many respects inclined to credit ourselves with too much, so here we perhaps
think too little of what we are capable of. We must here say with Aladdin:
``Whether I am master, a trial will show!'' Only by the experimental way can
we gain certainty here, as in so many other domains. If one should let
oneself be deterred from such trials by the thought that they are of no use,
since the necessary order of things goes its way in any case, one commits a
fallacy to which logic has even given a special name (``lazy reason,''
ignava ratio), and which would correspond to deriving from the lawfulness of
external nature that all our intervention in and action upon things was
fruitless. Precisely because our willing and acting are links in the series
of \opage{143}natural events can we modify them by altering the conditions
under which they occur.

Human nature is in itself neither good nor evil. Concepts such as good and
evil are relative concepts, which presuppose the application of a certain
point of view, a certain standard. If we start from the fact that human
nature is able to raise itself to a high degree of reason and freedom, of
truth and love, we shall rightly call it good. If, on the other hand, we
take account of the resistance that must be overcome and the struggle that
must be endured for this goal to be reached, one can call it evil. The
Christian doctrine of original sin and the Kantian doctrine of radical evil
are from this point of view fully justified; only a smooth, superficial
optimism can doubt it. The resistance that is offered, as well as the forces % misprint in the print: the point after "derom" wanting
that work forward, spring from the nature of humanity. J. G. \emph{Fichte}
sought to show that radical evil consists in sluggishness, inertia. Just as
in external nature there prevails a law of inertia, in that every thing
persists in its present state as long as it is not carried beyond it by
influences from without, so the human individual, without strong impulses
from without, will not develop to freedom and truth. Therefore life in
society, the intimate relation to the race, the whole historical side of
human life, are so decisive for ethical life. Our earlier investigations
have taught us from several sides to recognize the necessity of education.
Without the influence of \opage{144}authorities the spiritual growth of the
human race would not take place. The history of humanity is essentially the
history of authorities. In this great process of education the given, what
has been attained hitherto, which has come into equilibrium and stiffens
into fixed forms, offers resistance to the higher powers that are to work
their way forward, just as every body offers resistance to another that
would take its place. In the spiritual and especially in the ethical domain
the higher power must moreover gain entrance through man's freedom, and we
have seen that this consists in energetic attention and exertion; to this,
then, man must first be awakened: the powers he is to use he cannot give
himself, although he is to find them in his own nature. From sluggishness
there develop not only cowardice and falseness, as Fichte has shown, but
also defiance, the most intensive expression of inertia, which does not let
itself be carried beyond itself to take part in the life of a larger whole.
Defiance presupposes the less energetic forms of inertia, and only by here
again making the end the beginning can one regard it as the source of sin.
To overcome the resistance one must seek out the causes that originally
brought development onto this path; one must here, as in all education, find
motives and means whereby the isolation of the self can be abolished, its
interest being attached to an essential, universally valid end.

Humane ethics, for which the golden age or paradise is not a lost good but a
goal which, \opage{145}so far as is possible and justified under human
conditions, is to be reached by the labor of the race, cannot regard sin or
evil as something absolute. It knows a ``not yet,'' not a ``no more.'' It
knows that the goal must by its nature lie beyond the given, and that there
must therefore always be a distance between reality and ideal. This distance
it takes as given and does not try to fill it out with the help of fantasy.
Without such a distance, after all, it would not be worth living either.
(For a similar reason Baggesen asked always to keep seven unfulfilled wishes
in reserve). In so far as the distance is to be filled out in some other way
than by ever-renewed striving, this filling out is the business of poetry,
not of thought. But the distance must not be confused with irreconcilable
opposition; on the contrary, the imperfections themselves pave the way for
the more perfect. And the more evident it becomes that man really can make a
contribution to the formation of his own character, and thereby also of that
of the whole race, the more will interest also concentrate on this point,
and the more detours and sufferings will be spared. An eminent English
alienist, \emph{Henry Maudsley}, has recently insisted on the importance of
attention and interest for the development of character as a means of
preserving from insanity. He sees in the weakening of character, in the
general inclination to let oneself be driven by the current instead of
intervening self-actively, in great part the cause of the spread of
insanity. But he stresses that a \opage{146}result is reached here slowly
and perhaps first becomes visible when the work is continued through a
series of generations\footnote[1]{See the fine concluding chapter of his
book ``Le crime et la folie''. (Chap. IX: des moyens de se préserver de la
folie).}. Insanity stands in close connection with the other evils under
which humanity suffers; it is in itself only one of the forms under which
the barriers to higher development appear. The rational doctrine of the
freedom of the will which we have here attempted to maintain shows the
possibility of moving these barriers, even if in the single case it were
only by a hair's breadth. In ethics, as in other domains, we must often
reckon with infinitely small quantities. It at the same time insists that
this is possible only through penetrating psychological knowledge of
character and its conditions, just as we also master external nature only by
obedience to its laws. On the other hand, the assumption of an absolute
freedom would easily lead to fantastic attempts to recreate oneself, to
despair at not being able to, or to doubt of all lawful connection in the
spiritual domain altogether.

\begin{center}\rule{6em}{0.4pt}\end{center}

\end{document}
